\documentclass{mimosis}

\usepackage{fmtcount}
\usepackage{metalogo}
\usepackage{pifont}
\usepackage{adforn}
\usepackage{algpseudocode}
\usepackage{multicol}
\usepackage{wrapfig}

%%%%%%%%%%%%%%%%%%%%%%%%%%%%%%%%%%%%%%%%%%%%%%%%%%%%%%%%%%%%%%%%%%%%%%%%
% Some of my favourite personal adjustments
%%%%%%%%%%%%%%%%%%%%%%%%%%%%%%%%%%%%%%%%%%%%%%%%%%%%%%%%%%%%%%%%%%%%%%%%
%
% These are the adjustments that I consider necessary for typesetting
% a nice thesis. However, they are *not* included in the template, as
% I do not want to force you to use them.

% This ensures that I am able to typeset bold font in table while still aligning the numbers
% correctly.
\usepackage{etoolbox}
\usepackage{blindtext}

%%%%%%%%%%%%%%%%%%%%%%%%%%%%%%%%%%%%%%%%%%%%%%%%%%%%%%%%%%%%%%%%%%%%%%%%
% Hyperlinks & bookmarks
%%%%%%%%%%%%%%%%%%%%%%%%%%%%%%%%%%%%%%%%%%%%%%%%%%%%%%%%%%%%%%%%%%%%%%%%

\usepackage[%
  colorlinks = true,
  citecolor  = cardinal,
  linkcolor  = cardinal,
  urlcolor   = cardinal,
  % hidelinks  = false,
  unicode,
  ]{hyperref}

\usepackage{bookmark}

%%%%%%%%%%%%%%%%%%%%%%%%%%%%%%%%%%%%%%%%%%%%%%%%%%%%%%%%%%%%%%%%%%%%%%%%
% Bibliography
%%%%%%%%%%%%%%%%%%%%%%%%%%%%%%%%%%%%%%%%%%%%%%%%%%%%%%%%%%%%%%%%%%%%%%%%
%
% I like the bibliography to be extremely plain, showing only a numeric
% identifier and citing everything in simple brackets. The first names,
% if present, will be initialized. DOIs and URLs will be preserved.

\usepackage[%
  date = year,
  alldates=year,
  autocite     = plain,
  backend      = biber,
  doi          = true,
  url          = false,
  giveninits   = true,
  hyperref     = true,
  maxbibnames  = 3,
  maxcitenames = 1,
  sortcites    = true,
  style        = authoryear,
  uniquelist   = false,
  uniquename   = false,
  % sorting      = ynt,
  ]{biblatex}

\AtEveryBibitem{
    \clearfield{urlyear}
    \clearfield{urlmonth}
}

\input{bibliography-mimosis}
\addbibresource{Thesis.bib}
\setlength\bibitemsep{0.5\baselineskip}

%%%%%%%%%%%%%%%%%%%%%%%%%%%%%%%%%%%%%%%%%%%%%%%%%%%%%%%%%%%%%%%%%%%%%%%%
% Fonts
%%%%%%%%%%%%%%%%%%%%%%%%%%%%%%%%%%%%%%%%%%%%%%%%%%%%%%%%%%%%%%%%%%%%%%%%
\ifxetexorluatex
  \usepackage{unicode-math}
  \usepackage[lf]{ebgaramond}

  % Use sans serif font for headings only
  \setsansfont{Noto Sans Bold}
  \addtokomafont{chapter}{\sffamily}
  \addtokomafont{section}{\sffamily}
  \addtokomafont{subsection}{\sffamily}
  \addtokomafont{subsubsection}{\sffamily}

  % Set mono font for code
  \setmonofont[Scale=MatchLowercase]{JetBrainsMono Nerd Font Mono}

  % Load some missing symbols from another font.
  %
  % \setmathfont{STIXTwoMath}[%
  %   Path=./fonts/stix2-otf/,
  %   Extension = .otf,
  %   UprightFont = *-Regular,
  % ]

\else
  \usepackage[lf]{ebgaramond}
  \usepa
  \usepackage[oldstyle,scale=0.7]{sourcecodepro}
  \singlespacing
\fi

\usepackage{listings} % for code listings
\definecolor{codegreen}{rgb}{0,0.6,0}
\definecolor{codegray}{rgb}{0.5,0.5,0.5}
\definecolor{codepurple}{rgb}{0.58,0,0.82}
% \definecolor{backcolour}{rgb}{0.95,0.95,0.92}
% \definecolor{backcolour{rgb}{1.0,1.0,1.0}

\lstdefinestyle{mystyle}{
    backgroundcolor=\color{white},
    commentstyle=\color{codegreen},
    keywordstyle=\color{magenta},
    numberstyle=\tiny\color{codegray}\ttfamily,
    stringstyle=\color{codepurple},
    basicstyle=\linespread{1}\ttfamily\footnotesize,
    breakatwhitespace=false,
    breaklines=true,
    captionpos=b,
    keepspaces=true,
    numbers=left,
    numbersep=10pt,
    showspaces=false,
    showstringspaces=false,
    showtabs=false,
    tabsize=2,
    xleftmargin=15pt,
}
\lstset{style=mystyle}

\setlength{\parindent}{0pt}
\setlength{\parskip}{1em}

%%%%%%%%%%%%%%%%%%%%%%%%%%%%%%%%%%%%%%%%%%%%%%%%%%%%%%%%%%%%%%%%%%%%%%%%
% Acronyms and glossary
%%%%%%%%%%%%%%%%%%%%%%%%%%%%%%%%%%%%%%%%%%%%%%%%%%%%%%%%%%%%%%%%%%%%%%%%

\newacronym[description={Principal component analysis, an algorithm used for linear dimensionality reduction.}]{PCA}{PCA}{principal component analysis}

\newacronym[description={Gaussian mixture model, a clustering algorithm that fits mixtures of multidimensional Gaussian distributions.}]{GMM}{GMM}{gaussian mixture model}

\newacronym[description={Convolutional neural network, a neural network using convolutional feature extraction usually applied for 2D or 3D data.}]{CNN}{CNN}{convolutional neural network}

\newacronym[description={Multilayer Perceptron, a feed-forward neural network with fully connected hidden layers.}]{MLP}{MLP}{multilayer perceptron}

\newacronym[description={Electric organ discharge, the three dimensional electric dipole field generated by electric fish.}]{EOD}{EOD}{electric organ discharge}

\newacronym[description={Electric organ discharge frequency, how often a fish discharges its electric organ in one second.}]{EOD$f$}{EOD$f$}{electric organ discharge frequency}

\newacronym[description={Graphics processing unit, the unit of a computer that handles graphical rendering and consequently performs matrix operations very well, which is used for deep machine learning.}]{GPU}{GPU}{graphics processing unit}

\newacronym[description={Central processing unit, the main computing unit of a computer.}]{CPU}{CPU}{central processing unit}

\newacronym[description={Probability density function, describes the likelihood of a continuous random variable to take on a specific value within a given range.}]{PDF}{PDF}{probability density function}

\newacronym[description={Amplitude modulation, the variation in amplitude of a carrier signal.}]{AM}{AM}{amplitude modulation}

\newacronym[description={Short-term Fourier transform, a signal transform to compute spectrograms.}]{STFT}{STFT}{short-time Fourier transform}

\newacronym[description={Root mean square, a type of filter to estimate the envelope of a periodic signal.}]{RMS}{RMS}{root mean square}

\newacronym[description={Neural networks, algorithms inspired by the function and structure of the brain.}]{NN}{NN}{neural network}

\newacronym[description={Long short term memory network, a type of recurrent neural network with short term memory, commonly used for sequence data.}]{LSTM}{LSTM}{long short term memory}

\newacronym[description={Intersection over union, a metric to evaluate how strongly two bounding boxes overlap each other.}]{IoU}{IoU}{intersection over union}

\newacronym[description={Average precision, the area under the precision-recall curve, commonly used to evaluate the peformance of object detection algorithms.}]{AP}{AP}{average precision}

\newacronym[description={Mean average precision, the mean of the average precisions of all detection classes.}]{mAP}{mAP}{mean average precision}

\newacronym[description={Area under the curve, the definite integral of a curve.}]{AUC}{AUC}{area under the curve}

\newacronym[description={Receiver operating characteristic, a visual representation of the performance of a binary classifier.}]{ROC}{ROC}{receiver operating characteristic}

\newacronym[description={Generalized linear mixed effects models, a family of linear models capable of describing statistical relationships in nested and non-normal data.}]{GLMM}{GLMM}{generalized linear mixed effects models}

\newacronym[description={Variational autoencoder, a type of neural network that learns the probability distribution of a dataset and generates new samples based on this.}]{VAE}{VAE}{variational autoencoder}

\newacronym[description={Generative adversarial network, two neural networks. One network generates new data, while the other one evaluates how realistic the generated data is. This enables the model to learn to create novel samples.}]{GAN}{GAN}{generative adversarial network}

\newacronym[description={Generative pretrained transformer, a group of generative large language models using the transformer architecture.}]{GPT}{GPT}{generative pretrained transformer}

\newacronym[description={Bayesian information criterion, a metric used to fit models that rewards the model fit while penalizing the number of parameters used.}]{BIC}{BIC}{Bayesian information criterion}

\newacronym[description={Passive acoustic monitoring, biomonitoring using remote distributed microphones to record animal vocalisations.}]{PAM}{PAM}{passive acoustic monitoring}

\newacronym[description={Signal to noise ratio, compares the level of a signal of interest to that of the background noise.}]{SNR}{SNR}{signal to noise ratio}

\newacronym[description={Rectified linear unit, currently the most common activation function of artificial neural networks.}]{ReLU}{ReLU}{rectified linear unit}

% \newglossaryentry{LaTeX}{%
%   name        = {\LaTeX},
%   description = {A document preparation system},
%   sort        = {LaTeX},
% }
%
% \newglossaryentry{Real numbers}{%
%   name        = {$\real$},
%   description = {The set of real numbers},
%   sort        = {Real numbers},
% }

\makeindex
\makeglossaries

%%%%%%%%%%%%%%%%%%%%%%%%%%%%%%%%%%%%%%%%%%%%%%%%%%%%%%%%%%%%%%%%%%%%%%%%
% Ordinals
%%%%%%%%%%%%%%%%%%%%%%%%%%%%%%%%%%%%%%%%%%%%%%%%%%%%%%%%%%%%%%%%%%%%%%%%

\makeatletter
\@ifundefined{st}{%
  \newcommand{\st}{\textsuperscript{\textup{st}}\xspace}
}{}
\@ifundefined{rd}{%
  \newcommand{\rd}{\textsuperscript{\textup{rd}}\xspace}
}{}
\@ifundefined{nd}{%
  \newcommand{\nd}{\textsuperscript{\textup{nd}}\xspace}
}{}
\makeatother

\renewcommand{\th}{\textsuperscript{\textup{th}}\xspace}

%%%%%%%%%%%%%%%%%%%%%%%%%%%%%%%%%%%%%%%%%%%%%%%%%%%%%%%%%%%%%%%%%%%%%%%%
% Experiments
%%%%%%%%%%%%%%%%%%%%%%%%%%%%%%%%%%%%%%%%%%%%%%%%%%%%%%%%%%%%%%%%%%%%%%%%
% subsubsection (i think dont count bejond this level)
\setcounter{secnumdepth}{2}

% \KOMAoptions{DIV=last}
\definecolor{kelly_green}         {RGB}{ 76, 187, 23}
\definecolor{international_orange}{RGB}{255,  79,  0}
\definecolor{venetian_red}        {RGB}{200,   8, 21}
\definecolor{amber}     {RGB}{255,191,  0}
\definecolor{burgundy}  {RGB}{144,  0, 32}
% \definecolor{cardinal}  {RGB}{196, 30, 58}
\definecolor{red}  {HTML}{A2314E}
\definecolor{cardinal} {HTML}{00509d}
\definecolor{deep-azure}{RGB}{ 10, 50, 94}
\definecolor{mahogany}  {RGB}{192, 64,  0}
\definecolor{malory-red}{RGB}{192, 46, 16}
\definecolor{sky}       {RGB}{ 62, 93,240}
\definecolor{teal}      {RGB}{  0,128,128}
\definecolor{tyrian}    {RGB}{102,  2, 60}
\definecolor{yale}      {RGB}{ 70,130,180}
\definecolor{quality-high}  {RGB}{ 70,130,180}
\definecolor{quality-medium}{RGB}{255,191,  0}
\definecolor{quality-low}   {RGB}{196, 30, 58}


% Color for section numbers
% NOTE: Uncomment to enable colored section numbers
%
\renewcommand*{\chapterformat}{%
  \fontsize{60}{55}\selectfont\color{cardinal}\thechapter\autodot\enskip
}
\renewcommand*{\sectionformat}{%
  \textcolor{cardinal}{\thesection}\autodot\enskip%
}
\renewcommand*{\subsectionformat}{%
  \textcolor{cardinal}{\thesubsection}\autodot\enskip%
}
\renewcommand*{\subsubsectionformat}{%
  \textcolor{cardinal}{\thesubsubsection}\autodot\enskip%
}

% Image caption spacing
\setlength{\abovecaptionskip}{15pt plus 3pt minus 2pt}

%ornamet
\newcommand{\ornament}{%
  \begin{center}
    \vspace{-0.2cm}
    \large\textcolor{cardinal}{\ding{70}}
    \vspace{-0.2cm}
  \end{center}
}

%%%%% notes %%%%%%%%%%%%%%%%%%%%%%%%%%%%%%%%%%%%%%%%%%%%%%%%%%%%%%%%%%%%%%%%%%%
\newcommand{\note}[2][]{\textcolor{red}{[#1: #2]}}
\newcommand{\notepw}[1]{\note[PW]{#1}}

%%%%% incipit %%%%%%%%%%%%%%%%%%%%%%%%%%%%%%%%%%%%%%%%%%%%%%%%%%%%%%%%%%%%%%%%%
\title{Detection and Assignment of Electrocommunication Signals with Deep~Learning}
\subtitle{Agonistic Signaling in Weakly Electric Fish}
% \title{Automated Annotation and Emitter Recognition of Electric 
% Communication Signals}
% \subtitle{Deep Learning Applied to Animal Behavior}
% \title{Automated Detection of Communication Signals in Electric FishV}
% \subtitle{Deep Learning Applied to Animal Behavior}
% \title{Exploring Electrocommunication with Deep Learning}
% \subtitle{Agonistic Signaling in Weakly Electric Fish}

\author{Patrick Weygoldt}

\begin{document}

\frontmatter
\begin{titlepage}
  \vspace*{5cm}
  \makeatletter
  \begin{center}
    % \hrule
    {
      \setstretch{2}
      \begin{Huge}
        \@title
      \end{Huge}\\[0.5cm]
    }
    %
    \begin{Large}
      \@subtitle
    \end{Large}\\[1cm]
    %
    % \hrule
    % \begin{large}
    %   \adforn{49}
    % \end{large}\\[1cm]
    % %
    \emph{Masterarbeit}\\[1cm]
    %
    der Mathematisch-Naturwissenschaftlichen Fakultät der \\
    \textsc{Eberhard Karls Universität Tübingen} \\[1cm]
    %
    \emph{vorgelegt von}\\[0.5cm]
    \begin{large}
      \@author
    \end{large}
    %
    \vfill
    Tübingen, 29. Mai 2024
  \end{center}
  \makeatother
\end{titlepage}

\newpage
\null
\thispagestyle{empty}
\newpage

\begin{center}
  \textsc{Eigenständigkeitserklärung}
\end{center}
%
\noindent
%
Hiermit erkläre ich,

\begin{itemize}
  \item dass ich diese Arbeit selbst verfasst habe.
  \item dass ich keine anderen als die angegebenen Quellen benutzt und dass ich
  alle wörtlich oder sinngemäß aus anderen Werken übernommenen Aussagen als
  solche gekennzeichnet habe.
  \item dass die eingereichte Arbeit weder vollständig noch in wesentlichen 
  Teilen Gegenstand eines anderen Prüfungsverfahrens gewesen ist.
\end{itemize}
%
Tübingen, den 29. Mai 2024


\setcounter{tocdepth}{1}
\tableofcontents

\newpage

\begin{center}
  \textsc{Zusammenfassung}
\end{center}

\noindent

Kommunikation ist entscheidend im Wettbewerb um Ressourcen bei
Tieren, formt Dominanzhierarchien und kann die Notwendigkeit physischer
Auseinandersetzungen reduzieren. Die Funktion von \emph{chirps} – einer verbreiteten
Art elektrokommunikativer Signale bei elektrischen Fischen – ist 
umstritten. Bisherige Studien konzentrierten sich aufgrund der Schwierigkeiten
bei der Detektion von \emph{chirps} in freien Interaktionen vor allem auf isolierte
Fische. Diese Arbeit präsentiert eine \emph{Deep-Learning Pipeline} zur
automatisierten Erkennung von \emph{chirps} bei frei interagierenden
\textit{Apteronotus leptorhynchus}, um die Rolle von \emph{chirps} im ritualisierten
Wettbewerb um Unterschlupf und darüber hinaus zu erforschen.

Die entwickelte Pipeline nutzt \emph{Convolutional Neural Networks} für die
\emph{chirp}-Erkennung. Diese werden mit Spektrogrammen trainiert, auf denen
die \emph{chirps} markiert sind. Nach der Erkennung kommen speziell entwickelte
Algorithmen zum Einsatz, um die detektierten \emph{chirps} dem jeweils
kommunizierenden Fisch zuzuordnen. Diese Methode ermöglicht eine präzise
Identifizierung und Zuordnung von \emph{chirps} zu einzelnen Fischen und führte
zu einem Datensatz von über 73.000 \emph{chirps}, die bei freien Interaktionen
zwischen zwei Fischen produziert wurden.

In ersten Analysen wurden die Muster der \emph{chirps} und deren
Verhaltenskorrelationen untersucht. Eine wichtige Erkenntnis war die starke
Korrelation zwischen \emph{chirps} von submissiven Individuen und dem Abbruch
aggressiver Interaktionen bei Wettbewerben um Unterschlupf. Diese Beobachtung
legt nahe, dass \emph{chirps} über die sensorische Wahrnehmung hinaus eine
bedeutende Rolle spielen. Die von submissiven Fischen ausgesendeten
\emph{chirps} schienen jedoch keinen messbaren Einfluss auf das Verhalten der
dominanten Fische zu haben, was die Notwendigkeit weiterer Untersuchungen
hervorhebt.

Darüber hinaus lässt sich diese automatisierte Methode auch auf Felddaten
anwenden. Vorherige Studien haben gezeigt, dass die markerlose Verfolgung
mehrerer Individuen in ihrer natürlichen Umgebung über die Erfassung
elektrischer Felder auf Elektroden-Gittern möglich ist. Die \emph{Chirp
Detection Pipeline} baut auf diesen Erkenntnissen auf und eröffnet die
Möglichkeit, das gesamte elektrokommunikative Repertoire aus solchen Aufnahmen
zu extrahieren. Diese Entwicklung bietet erstmals die Gelegenheit, umfassende
Informationen über Identität, Geschlecht, Bewegungen und nun auch
Kommunikationssignale für quantitative ethologische Studien von ungestörten
Populationen anhand von Daten aus einer einzigen Modalität zu analysieren.

\newpage

\begin{center}
  \textsc{Abstract}
\end{center}

\noindent

Effective communication plays a crucial role in resource competition among
animals, shaping dominance hierarchies and reducing the necessity for physical
confrontations. The function of chirps, a common form of
electrocommunication signal in electric fish, is still a topic of debate. Past
studies have mainly focused on isolated fish due to the technical limitations
in detecting chirps during interactions. This thesis introduces a deep learning
pipeline for automated chirp detection in freely moving Brown Ghost Knifefish
to investigate the function of chirps in staged competition scenarios and
beyond.

The pipeline applies deep convolutional neural networks for chirp detection,
which are fine-tuned using labeled spectrograms of chirping fish. Following
detection, it uses custom algorithms to link detected chirps with the specific
fish producing them. This methodology facilitated the identification and
allocation of chirps to individual fish, resulting in a substantial dataset of
over 73,000 chirps detected during unconstrained interactions.

Initial analyses focused on the patterns of chirping and their behavioral
associations. A notable finding is a significant correlation between chirps
emitted by subordinate individuals and the cessation of aggressive interactions
when competing for shelter. This observation strongly implies that chirps serve
a purpose beyond active sensing. Still, the chirps emitted by the subordinate
fish did not seem to influence the behavior of the dominant fish in a
measurable manner, which highlights the need for further investigation.

Beyond laboratory settings, this automated methodology can be applied to field
recordings as well. Prior research has illustrated the feasibility of
markerless tracking of multiple individuals in their natural environments by
capturing their electric fields on electrode grids. Building on this
foundation, this approach provides the capability to extract the complete
electrocommunication repertoire from such recordings. This advancement now
represents the opportunity to extract information concerning identity, sex,
movements, and communication signals for quantitative ethological studies in
undisturbed populations using data from a single modality for the first time.

\vspace*{5cm}
\makeatletter
\begin{center}
  Contempt for simple observation is a lethal trait in any science.\\[0.5cm]
  — \textsc{Niko Tinbergen} in \enquote{On Aims and Methods of Ethology}\\
  \vfill
\end{center}
\makeatother

\begin{center}
  \textsc{Acknowledgements}
\end{center}

\noindent

I would like to express my deepest gratitude to Jan Benda, my thesis advisor,
for his support and guidance throughout the research and writing process of
this thesis. Thank you for introducing me to both programming and electric
fish, and for the opportunity to observe them in their natural habitat in
Brazil. Thank you also for entrusting me with this project and the freedom to
experiment with and venture into the field of deep learning.

I am also thankful to Nico Michiels for the many unique opportunities to learn
and grow during my studies. Thank you for the significant boost to my diving
skills, achievements that would not have been possible without my stays in
Indonesia and Corsica, which you facilitated. I am confident that these skills
will be invaluable during future field trips.

My thanks also go to Till Raab and Jan Grewe at the Neuroethology Group, whose
expertise and knowledge were crucial during my studies. Thank you, Till, for
introducing me to the analysis of electrode grid datasets and for supervising
the initial version of the chirp detector. And thank you, Jan, for always being
ready to resolve urgent programming questions, regardless of the time or day.

It was truly a pleasure working with all of you.

On a personal note, I must acknowledge the generous financial support from my
beloved family. My heartfelt gratitude goes to my mother, father, and maternal
grandmother, whose dedication and sacrifices have enabled me to pursue my
academic goals. Their belief in my potential has been a constant source of
motivation throughout this journey. Additionally, I want to thank my maternal
grandfather, who ignited my passion for nature and animals at an early age.

Among my greatest supporters has been my partner, Sofie. You brought me to
Tübingen and have been my steadfast companion throughout all of my studies.
Your unending emotional support, patience, and love have been my rock during
all academic challenges. I am deeply grateful for the countless days you spent
quizzing me with flashcards, and for your genuine interest in and thorough
understanding of my research, despite your different academic background. Our
walks in the forests around Tübingen sparked many new ideas. Thank you for
always being there, especially through the tougher times.

Lastly, I extend my gratitude to all those who were directly or indirectly
involved in the successful completion of this thesis. Your contributions have
been invaluable.


\mainmatter

\chapter{Introduction}
\label{chap:intro}

Communication is a fundamental aspect of life on earth, essential for conveying
information, coordinating behaviors, and establishing social relationships.
From intricate visual displays to complex vocalizations, communication plays a
crucial role in the survival and reproduction of species across the tree of
life. In humans, communication takes on a myriad of forms, from spoken language
over gestures and facial expressions to written text. The study of animal
communication offers insights into the ecology and evolution of
species, providing a window into the underlying mechanisms and functions of
different communication signals used in various contexts. At the same time, it
provides a platform that could help us understand how communication generally
and our language specifically evolved and how it is produced and perceived in
the brain.

\ornament

Electric fish offer a unique opportunity to study animal communication and
social behavior but are rarely used as a model system outside the field of
Neuroethology. Beyond utilizing their electric field for active sensing,
wave-type electric fish have been used extensively to study the physiological
underpinnings of the production and perception of communication signals. The
most frequently observed, and hence studied, electrocommunicaion signal are
so-called chirps. During a chirp, electric fish briefly and transiently
increases the frequency with which it discharges its electric organ, a
phenomenon found in many species
\parencite{petzoldCoadaptationElectricOrgan2016,
smithEvolutionElectricCommunication2016,
freilerElectrocommunicationSignalsAggressive2022,
zhouStructureSexualDimorphism2006, hoSexDifferencesElectrocommunication2010}.

Understanding chirps necessitates a introduction of the \acrfull{EOD} in
wave-type weakly electric fish. The EOD, a quasi-sinusoidal waveform similar to
a constant-pitch tone in audio, features a fundamental frequency
\acrshort{EOD$f$} and harmonics. Its amplitude decays dependent on the fish's
orientation, body curvature, and distance relative to the recording electrode
\parencite{bendaPhysicsElectrosensoryWorlds2020}. While EOD$f$ differs among
individuals, it's generally stable, with variations mainly due to water
temperature changes \parencite{moortgatSubmicrosecondPacemakerPrecision1998}.
Frequencies range from less than \SI{50}{\hertz} to over \SI{2}{\kilo\hertz},
with each species occupying a characteristic frequency range
\parencite{bendaPhysicsElectrosensoryWorlds2020,
evansTaxonomicRevisionDeep2017a}. Despite its simplicity, the EOD waveform can
reveal species-specific traits under standardized recording conditions.

Wave-type weakly electric fish can modulate the frequency of their EOD on time
scales of milliseconds to minutes to produce various communication signals
\parencite{zakonEODModulationsBrown2002a}. In most cases, the chirps of
\textit{A.~leptorhynchus} are \enquote{simple} up-modulations of the EOD$f$
that last between \SI{20}{\milli\second} and \SI{500}{\milli\second} and can
reach between \SI{20}{\hertz} and \SI{500}{\hertz} above the baseline EOD$f$.
From now on, when I refer to a chirp as \enquote{short} or \enquote{long}, I
will refer to the duration of the chirp. When mentioning the properties
\enquote{small} or \enquote{large}, I will refer to the frequency up-modulation
from baseline. Chirps can come in many shapes: Small and short chirps are
usually approximately shaped like a Gaussian as illustrated in figure
\ref{fig:multifish_chirps} A. If they become larger, this Gaussian can be
followed by a brief undershoot. In contrast, the shape of long and large chirps
usually vaguely resembles a square wave as in figure \ref{fig:multifish_chirps}
E. If we now take a look at the signal that underlies this change in frequency,
so the signal that we could record, we can see, that the third feature of a
chirp, in addition to the duration and the frequency change, is the change in
amplitude. For small chirps (figure \ref{fig:multifish_chirps} B) this trough
in amplitude is usually not as pronounced as it is for larger chirps (figure
\ref{fig:multifish_chirps} F).

\begin{figure}
  \centering
  \includegraphics[width=\textwidth]{../figures/multifish_chirps.pdf}
  \caption{
    Overview of how chirps as changes in the EOD$f$ of a fish can be recorded
    in different settings if other fish are around, simulated from the
    instantaneous frequency evolutions from chirps of real fish. Chirps are
    defined as a brief elevation of the EOD$f$ of one fish. They can be
    relatively small and short (A) or long with a high up-modulation (E). If
    only the chirping fish would be recorded in a chirp-chamber-like setup, a
    small chirp is already hardly visible by just the EOD alone (B) while a
    larger chirp might appear a bit clearer (F). In the field however, the EODs
    of multiple fish superimpose, which makes chirp detection but also chirp
    assignment much harder without a frequency spectrum (C, G). Because of
    this, we need a \acrfull{STFT} to convert the recordings to spectrograms,
    which allows an approximation of the underlying discharge frequencies (D,
    H). As this is just an approximation, the resolvability of chirps and the
    separability of fish forms a trade-off: The temporal resolution of the
    spectrogram, if sufficient to separate fish, is not sufficient to resolve a
    full chirp (D) but rather produces an \enquote{artifact} caused by the
    presence of the chirp which is usually much longer than the underlying
    chirp, particularly for small chirps. For long chirps (H), this is not so
    much of a problem.
  }
  \label{fig:multifish_chirps}
  \vspace{0.5cm}
\end{figure}

Chirps offer a unique opportunity to link advances in animal behavior to the
underlying neural mechanisms: They are extensively studied in controlled
conditions and the endocrinological as well as neurological mechanisms to
produce and perceive them are comparatively well understood
\parencite{hupeEffectDifferenceFrequency2008,
zupancElectrogenesisElectroreceptionOverview2005,
smithSerotonergicActivation5HT1A2008, dunlapGlucocorticoidReceptorBlockade2011,
dulkaTestosteroneModulatesFemale1994, kellerControlPacemakerModulations1991,
kraheNeuralMapsElectrosensory2014, marsatCellularCircuitProperties2012,
dunlapVocalElectricFish2021, bendaSpikeFrequencyAdaptationSeparates2005,
bendaSynchronizationDesynchronizationCodeNatural2006}. But the role of chirps
produced in various naturalistic behavioral contexts is still mysterious.
Functionally, chirps as a communication signal have been attributed to roles in
courtship as well as signaling dominance and submission, the latter two seeming
contradictory \parencite{batistaNonsexbiasedDominanceSexually2012,
hagedornCourtSparkElectric1985a, englerDifferentialProductionChirping2001,
henningerStatisticsNaturalCommunication2018}. On the other hand, recent work
suggested that chirps might not serve as a means of conventional communication
at all but as a form of autocommunication to increase the capabilities of the
active electric sense in social contexts \parencite{obotiWhyBrownGhost2023}.
Except for the work by \textcite{henningerStatisticsNaturalCommunication2018},
most work underlying proposals for various functional roles of chirps rely
heavily on data recorded in artificial conditions and none of them provide
causal evidence. But to correctly interpret physiological results holistically,
we must include a species ecology and evolution
\parencite{tinbergenAimsMethodsEthology1963}. A dramatic example is the work by
\textcite{henningerStatisticsNaturalCommunication2018}, who showed that the
stimulus regimes commonly used in controlled conditions to study the neural
mechanisms underlying electric communication are not representative of the
natural environment. I argue that before probing for potential explanations for
a specific behavior in highly controlled conditions, we should aim to first
\emph{observe} correlations that suggest causal relationships to generate
\emph{informed} hypotheses before we test them with controlled experiments. In
other words: We need to observe which fish chirps at what time and during which
kind of social interactions in a naturalistic context. 

The biggest challenge preventing these approaches to studying chirping is that
chirps are hard to detect automatically in behaviorally relevant settings.
Currently, chirps can be automatically detected in cases where the fish are
separated from each other so that their electric fields do not superimpose
\parencite{lehotzkySupervisedLearningAlgorithm2023,
eskeEffectUrethaneMS2222023, fieldJARChirpsGymnotiform2019}. But studies where
fish are kept together are much more informative about their natural behavior
and often uncover behaviors that challenge common assumptions from previous
literature \parencite{henningerStatisticsNaturalCommunication2018,
fortuneSpookyInteractionDistance2020a}. In these cases, the electric fields of
the fish superimpose and the chirps are much harder to detect
\parencite{obotiWhyBrownGhost2023,henningerStatisticsNaturalCommunication2018}.
This is either not studied at all, or if so involves a lot of manual
annotation, which is a labor-intensive and error-prone process and not feasible
for large datasets, such as the ones recorded in the field. But it is not the
problem that we do not have the tools to detect chirps in these cases. As an
example, smartphone apps can automatically recognize bird songs and classify
the species just from comparatively poor-quality phone recordings
\parencite{kahlBirdNETDeepLearning2021}. Despite this, many studies still rely
on manually labeling communication signals of all kinds of species. Also,
software such as
\href{http://www.mackenziemathislab.org/deeplabcut}{\texttt{DeepLabCut}} and
\href{https://sleap.ai/}{\texttt{SLEAP}} enables automated tracking of
individual body parts of most animals \parencite{pereiraSLEAPDeepLearning2022a,
lauerMultianimalPoseEstimation2022}. This provides the possibility of high
throughput in behavioral experiments and frees time that can be allocated to,
for example, recording more data. Algorithms to automate the conversion of
unstructured data (such as video, audio or electric recordings) into structured
data (such as tables to do statistics on) are widely available, they just need
to be embedded into software packages to make them accessible to researchers.

In this thesis, I want to provide a demonstration of such a fully automated
pipeline. Here, I aim to implement a first proof of concept for fully automated
chirp detection in large datasets recorded in the field for \textit{A.
leptorhynchus}, although the pipeline can be extended to any chirping species
of wave-type electric fish. As you are about to see, this does not require
inventing new technologies but simply connecting existing, openly available and
well documented tools into a data pipeline to fully automate the detection of
chirps in large datasets for quantitative research in electrocommunication.
This will not only enable to observe and hence analyze chirps in field
recordings but could also be an important building block for electric fish
monitoring.

To understand the features that are useful to detect chirps, I will first
introduce the physical properties of the \acrfull{EOD} of weakly electric and
how chirps change it in the \hyperref[chap:background]{next} chapter. We will
then transition to a brief introduction on neural networks and how they can
applied to the problem of chirp detection. In the
\hyperref[chap:methods]{following} chapter, I will introduce the training data
generation pipelines for machine learning models, how the models are trained
and how their performance is evaluated. I will also explain how the detector
can be applied to a large dataset of electric recordings and how I analyzed
behaviors associated with chirps. In the \hyperref[chap:results]{results}, we
will look at the full architecture of the chirp detection pipeline, the
performance of the neural networks and the results of the analysis of the large
dataset of weakly electric fish recordings. And finally, in the
\hyperref[chap:discussion]{discussion} chapter, we will discuss the
implications of the results and the limitations of the pipeline. We will also
discuss how the pipeline can be improved and how it can be applied to other
species and other tasks beyond analyzing animal communication, such as
environmental monitoring. For now, let's start with the fundamentals: How do we
record the electric fields of weakly electric fish?


\chapter{Technical background}
\label{chap:background}

Detection tasks can be solved with a variety of solutions, each tailored to
specific problems and desired performance. The detection of chirps is no
exception. In this chapter, we will introduce some necessary background
information: How are fields of electric fish recorded and how does this relate
to chirp detection? What can we observe chirps if fish can freely interact? How
were chirps detected before? And what is an artificial neural network?

\section{How to record electric fish?}
\label{sec:recording}

Traditionally, electric fish recordings have utilized a differential setup with
two electrodes, one placed near the head and the other near the tail, to
capture the electric field generated by the fish. This setup, known as a
\enquote{chirp chamber}, is widely used in research (e.g.,
\cite{dunlapProductionAggressiveElectrocommunication2003,
fugereElectricSignalsSpecies2010}). Species like \textit{A.~leptorhynchus}
often remain in shelters when exposed to light, mimicking their natural
inactive phase, which reduces the need for physical restraint during
recordings. This method produces a nearly perfect, quasi-sinusoidal
\acrfull{EOD} from single, isolated individuals. Such a recording is ideally
centered around zero and noiseless enough so that every single period, i.e.,
every single discharge as well as the species specific waveform, is clearly
visible. However, it limits the study of social behaviors and interactions, as
animals cannot engage physically or electrically, making it less suitable for
research into communication.

To investigate such interactions, the traditional setup can be adapted by
housing two fish in a single tank divided by a mesh, allowing electrical but
not physical interaction (e.g., \cite{obotiWhyBrownGhost2023}). Electrodes are
placed in each compartment to record the electric fields. While these
recordings can capture the fish's instantaneous frequencies, they tend to be
noisier and are better analyzed using spectrograms rather than direct frequency
calculations. However, such setups are highly dependent on the experimental
conditions, including tank size and electrode placement, which can influence
the observed behaviors and may not accurately reflect natural interactions.
Smaller setups are often used to minimize electrode requirements, but this can
lead to unnatural behaviors not observable in the wild. Additionally, the
effectiveness of EOD recording varies with electrode spacing, and physical
barriers limit genuine interactions. A potential solution is to employ an
electrode grid, distributing many electrodes evenly throughout the setup to
enhance recording fidelity during interactions
\parencite{junLongtermBehavioralTracking2014, pedrajaUseSupervisedLearning2021,
henningerTrackingActivityPatterns2020,
madhavHighresolutionBehavioralMapping2018}.

Electrode grids, initially developed for field research, are equally effective
in controlled studies, where they ensure complete coverage of the recording
setup and allow free interactions between fish. These setups necessitate the
use of spectrograms for frequency analysis, which introduces challenges related
to spectrogram parameters and tracking individual fish EODs over time.
Solutions to these challenges have been proposed by
\textcite{henningerTrackingActivityPatterns2020} and
\textcite{madhavHighresolutionBehavioralMapping2018}, and further refined by
\textcite{raabAdvancesNoninvasiveTracking2022}, culminating in the
\href{https://github.com/tillraab/wavetracker}{\texttt{wavetracker}} software
package for automatic fundamental frequency tracking. This advancement
facilitated the first study on electrocommunication in unconstrained
interactions in \textit{A.~leptorhynchus} on a large dataset by
\textcite{raabElectrocommunicationSignalsIndicate2021}. However, chirp
detection still requires manual spectrogram analysis for fast frequency
changes.

As we transition to the next section, I will cover previous attempts to detect
chirps from the rich datasets generated by electrode grids. This exploration
will not only mark the challenges faced in earlier methodologies but allow a
discussion on the techniques that have emerged to overcome these obstacles. But
to understand the challenge of detecting chirps in the EOD of multiple fish, it
is crucial to first introduce how we can make chirps accessible in recordings
of multiple fish.

\section{How to observe chirps in naturalistic settings?} % (fold)
\label{sec:chirps}

In the \hyperref[chap:intro]{introduction} we already established that the main
features of the \acrshort{EOD} that are changed during a chirp are the
frequency of discharge as well as the amplitude. But this amplitude feature is
not easily accessible when multiple fish are recorded on the same electrode,
which usually happens when fish interact. In this case, the resulting signal is
a superposition of the EODs of all fish, where the amplitude is modulated by a
so-called \emph{beat}. The exact shape of this \acrfull{AM} is dependent on the
differences between the discharge frequencies of the fish, the distance of each
fish to the recording electrode and even the orientation as well as curvature
of the body of single fish. Particularly small and short amplitude modulations
during chirps often drown in this complexity (figure \ref{fig:multifish_chirps}
C) while the AM of larger chirps remains visible in some cases (figure
\ref{fig:multifish_chirps} G). But in these cases, the time scales mix with
those of body movements. So if such a modulation is the result of a chirp or
that of a brief attenuation by a non-conducting object, such as a stone, is not
always clear.

So, in most species the property of the signal that firstly changes the
\enquote{most} and secondly is not influenced by external factors is the
frequency modulation during a chirp. To approximate how frequency components
change over time, we can use spectrograms. To compute a spectrogram, a
\acrfull{STFT} is applied to the recorded EOD to estimate the frequencies and
their powers of a part of the signal. The main pitfall with this approach is
given by the uncertainty principle: To get a very precise estimate of the
powers of certain frequencies in a signal, we need to compute the STFT over a
sufficiently large chunk of the signal, as the duration of a single segment
sets the frequency resolution. As an example, including three seconds of a
signal into a single Fourier transform results in a very fine frequency
resolution of $1/3$ \si{\hertz}. That means that even fish that are just one
Hertz apart in their EOD$f$ would be separable from each other on a
spectrogram. But for each point in time, we would end up with just the average
of the frequency estimate in three seconds around this point. Given that a
chirp only lasts fractions of a second, it would be impossible to find chirps
on these spectrograms. On the other hand, if we resolve the time domain of a
spectrogram to be able to closely follow the frequency modulation of a chirp,
the frequency bands for individual fish would be so broad that they would merge
into each other. Now we see chirps well, but cannot tell which fish produced
them. The ideal spectrogram considers the optimal trade-off between the two: It
has sufficient temporal resolution to make chirps somewhat visible but also
maximizes the separability of fish. An example for such a spectrogram is
illustrated in figure \ref{fig:multifish_chirps} D for small chirps and figure
\ref{fig:multifish_chirps} H for large chirps. We see that particularly the
temporal evolution of small chirps become blurry as the temporal resolution is
decreased in favor to make fish separable. To resolve chirps while still
maintaining reasonable separability of individual fish, a spectrogram that uses
bins of approximately \SI{20}{\milli\second} corresponding to 4096 samples for
a signal with a \SI{20}{\kilo\hertz} sampling rate is ideal. This makes sense
given that chirps are usually not shorter than \SI{20}{\milli\second} and at
least a single frequency estimate per chirp is needed to resolve it well. The
second important factor is the overlap between the bins: If we estimate the
frequency once every \SI{20}{\milli\second} and just by chance place one bin at
the start and one bin at the end of the chirp, the frequency in these bins will
still be elevated compared to the baseline but the increase won't be that
pronounced. So ideally, the overlap of the windows in which the frequencies are
estimated is large and traded off against the availability of computational
power. A stronger overlap produces a finer spectrogram but requires more
computations and memory.

Now that we have introduced why chirps are interesting, which parameters define
chirps in naturalistic recording conditions and how we can make them
approachable by computing specific types of spectrograms, let us discuss
previous implementations of chirp detectors.

\section{How were chirps detected in the past?}
\label{sec:chirpdetectors}

Detecting chirps in the EOD of electric fish, particularly in a controlled
environment with restrained fish, has traditionally been straightforward. In
such settings, fish are placed in chirp chambers where their movements are
limited, ensuring stable EOD signals without overlap from other sources. The
most common technique involves band-pass filtering the EOD signal and
calculating the instantaneous frequency through zero-crossings, with chirps
identified by rapid frequency increases as illustrated in figure
\ref{fig:multifish_chirps} A and E.

\textcite{henningerStatisticsNaturalCommunication2018} developed a method to
detect chirps in the electric organ discharges of multiple fish on field
electrode grid recordings. They initially filtered each fish's EOD within a
$\pm$ \SI{7}{\hertz} range around the baseline frequency, and additionally
within a \SI{14}{\hertz} window positioned \SI{10}{\hertz} above the baseline
frequency to capture the chirp signals, which momentarily elevate into this
range. To identify chirps, they employed a \acrfull{RMS} filter to estimate the
signal's instantaneous amplitude or envelope, looking for peaks within this
\enquote{search envelope} that exceeded a predefined threshold. This method
also involved monitoring for concurrent drops in the EOD amplitude, with
simultaneous amplitude drops and search envelope peaks indicating chirps. Each
detected chirp was manually verified and attributed to the emitting fish based
on the assumption that such signal patterns uniquely occur during chirping.
This approach enabled the successful detection and assignment of chirps to
specific fish in complex multi-fish recordings during courtship behavior.

In their research on the role of chirps in \textit{A. leptorhynchus},
\textcite{obotiWhyBrownGhost2023} applied a technique akin to
\textcite{henningerStatisticsNaturalCommunication2018} for chirp detection,
focusing on identifying peaks in the power spectral density between the
fundamental frequency and the first harmonic of the EOD. Peaks surpassing a
noise threshold were classified as chirps, with each detection undergoing
manual verification. For chirp assignment to specific fish, Oboti
leveraged the recording setup's electrode spacing, ensuring one fish dominated
the signal on a particular electrode, thus facilitating accurate chirp
attribution. While effective in controlled lab settings, this method's
requirement for limited fish movement compromises its applicability to field
recordings and may not fully capture natural fish communication behaviors due
to the imposed movement restrictions.

Neither study discussed the false positive rates of their chirp detection
methods or the extent of manual verification needed. In field recordings, it is
possible to capture up to 25 fish simultaneously over several weeks, with fish
potentially chirping at rates of two to three times per second, sometimes for
hours. Henninger's method struggles with chirp assignment in recordings where
fish EOD frequencies are closely spaced. Oboti's approach, reliant on
restricting fish movement, probably faces challenges in accurately assigning
chirps, since chirps often occur when fish are in close proximity, making it
difficult to assign chirps based on the strongest electrode signal.
Additionally, both methods are susceptible to interference from electrical
noise, such as white noise bursts that mimic chirps and can interfere with the
EOD signal, complicating chirp identification without extensive manual review,
which is impractical for large datasets.

The following sections will detail the advancements in chirp detection and
different approaches that I applied in the preparation and in initial
experiments of this thesis.

\subsection{Dynamic filtering of the EOD for chirp detection}
\label{sec:filterdetector}

Our initial chirp detection approach, developed with Sina Prause, Alexander
Wendt, and Till Raab, utilized dynamic filtering based on simple filters and
parameters from \textcite{henningerStatisticsNaturalCommunication2018}. Unlike
Henninger's stationary search envelope, our method dynamically adjusted the
frequency range based on the presence of other fish, optimizing chirp detection
for each recording window. We also incorporated spatial information, focusing
on the electrode with the strongest signal for each fish, to enhance
\acrfull{SNR} and avoid frequency overlap with other fish. The detection relied
on peaks occurring simultaneously in three features: The dynamic search
envelope, the envelope filtered around the EOD baseline, and the instantaneous
frequency of the baseline EOD (this code is available on GitHub:
\url{https://github.com/weygoldt/gp-chirpdetector}).

This method was effective for simple scenarios with two fish with a wide gap in
their EOD$f$s. However, it struggled with fast moving fish due to noisy
baseline envelopes and unreliable spatial information, leading to a high false
positive rate and difficulty in assigning chirps to the correct fish,
especially for long chirps. Given these limitations and the reliance on
hand tuned parameters, I shifted towards a machine learning approach to better
capture the nuanced features distinguishing chirps from noise, which will be
discussed in the following section.

\subsection{Sliding convolutional classifier for chirp detection in spectrograms}
\label{sec:cnndetector}

I next employed a sliding \acrfull{CNN} to identify chirps, leveraging their
visibility in spectrograms (as shown in figure \ref{fig:multifish_chirps}).
This method involved classifying spectrogram snippets as \enquote{containing a
chirp} or \enquote{not containing a chirp} using a CNN. That is a neural
network adept at processing images through convolutional layers that extract
relevant features for classification. The process entailed moving a sliding
window across the spectrogram to generate snippets, which the CNN then
evaluated by assigning a score indicating the likelihood of a chirp's presence.
Scores were low-pass filtered and thresholded to identify chirps. The CNN was
trained on simulated chirps—simple Gaussian modulations in the frequency domain
with added noise. This model was then applied to detect chirps in real
multi fish recordings (this code is available on GitHub at:
\url{https://github.com/weygoldt/cnn-chirpdetector}).

However, this approach had limitations as well. While this method improved upon
dynamic filtering in detecting small and short chirps, it failed to reliably
identify and assign large and long chirps, which are of particular interest due
to their sparse occurrence and potential significance in specific contexts.
Additionally, the CNN struggled to differentiate between large chirps and
certain noise artifacts, and assigning chirps to the correct fish was
unreliable due to the fixed spatial context of the sliding window. In
conclusion, the sliding CNN marked progress in comparison to the dynamic
filtering approach, but it fell short in automatically detecting chirps in
complex multi-fish recordings. The following sections will detail a more
sophisticated neural network pipeline developed to address these challenges.

\vfill

\section{What is an artificial neural network?}
\label{sec:nns}

Before introducing our final approach, it's crucial to grasp the fundamentals
of a \acrfull{NN}, a mathematical model inspired by the structure of brains.
NNs are a subset of supervised learning within machine learning, designed to
learn complex non-linear relationships between input and output data, such as
classifying images into categories like \enquote{cat} or \enquote{dog}.

A NN consists of interconnected nodes or neurons arranged in layers: An input
layer receiving data, hidden layers processing the data, and an output layer
delivering the model's prediction. For instance, to classify a $(500 \times
500)$ pixel image, the input layer would have $(500 \times 500 \times 3)$
neurons for each pixel and color channel, while the output layer might have two
neurons representing the categories \enquote{cat} and \enquote{dog}. Neurons
across layers are linked by weights, the learnable parameters of the model.
Training an NN involves adjusting these weights to minimize the difference
between the model's predictions and the actual labels, a process facilitated by
an algorithm called backpropagation. This algorithm calculates the gradient of
the loss function (a measure of prediction error) with respect to the weights
and updates the weights to reduce the loss.

The simplest model, the \acrfull{MLP}, comprises multiple fully connected
layers as illustrated in figure \ref{fig:mlp}. However, MLPs alone often
struggle with complex real-world data, either failing to learn relevant
features or overfitting. This challenge is addressed by incorporating
specialized layers for different data types, such as \acrfull{LSTM} blocks for
time-dependent data or \acrshort{CNN}s for images. CNNs, for example, use
two-dimensional convolutions to extract spatial features from images, mimicking
the human visual system's ability to recognize edges, shapes, and textures.

\begin{figure}[!ht]
  \centering
  \includegraphics[width=\textwidth]{../figures/mlp.pdf}
  \caption{
    Schematic representation of a MLP with an eight-dimensional input and a
    two-dimensional output with three hidden layers. Each node is a neuron that
    is connected to all other neurons in the next layer. The thickness of the
    connecting lines corresponds to the weights (that are randomly initialized
    here). As a result, this network could be used to classify an eight
    dimensional input into two classes. Visualization by \textcite{NNSVG}.
  }
  \label{fig:mlp}
  \vspace{0.5cm}
\end{figure}

These feature extraction layers, or the \enquote{backbone} of the network, are
crucial for learning data-specific features and are followed by a
\enquote{head} responsible for making predictions. The architecture of the head
varies with the task, ranging from a single neuron for binary classification to
a softmax layer for multi-class classification or a set of neurons for object
detection, i.e., regression tasks.

In summary, NNs are powerful tools for learning from and making predictions
about complex data. Their architecture, comprising feature extraction layers
and prediction heads, can be tailored to a wide range of tasks, from text
generation to object detection. With this foundation, we can now explore how
CNNs are applied to detect chirps in recordings of multiple fish, leveraging
their ability to extract and interpret spatial features from spectrograms. The
next section will introduce the application of CNNs for chirp detection in
multi-fish recordings, showcasing the practical use of NNs in analyzing complex
biological data.

% \section{What is an Artificial Neural Network?} % (fold)
%
% Now that we already introduced the concept of a neural network, before we
% proceed to the final approach, it is important to understand the basics of
% neural networks.
%
% A neural network is a mathematical model that is inspired by the human brain.
% Neural networks fall under the category of supervised learning into the
% broader field of machine learning, meaning that they can learn a complex
% non-linear transformation between input data, for example images, and output
% data, for example the names of animals that might be in those images. It is a
% collection of connected nodes, called neurons, that are organized in layers.
% The first layer is the input layer, the last layer is the output layer and
% all layers in between are called hidden layers. The input layer is the layer
% that receives the input data, the output layer is the layer that returns the
% output of the model and the hidden layers are the layers that process the
% input data to return the output. As an example, if you would want to classify
% an image of 500 by 500 pixels with three color channels into two classes,
% let's say a \enquote{cat} or a \enquote{dog}, the input layer would consist
% of $(500 \times 500 \times 3)$ neurons and the output layer would consist of
% two, one for \enquote{cat} and one for \enquote{dog}. Classes are usually
% encoded by integer numbers, lets say 0 for \enpquote{dog} and 1 for
% \enquote{cat}.
%
% The neurons in each layer are connected to the neurons in the next layer. The
% connections between the neurons are called weights. The weights are the
% parameters of the model that are learned during training. The model is
% trained by adjusting the weights so that the output of the model is as close
% as possible to the true output. So neuron could be expressed as a function
% that takes the weighted sum of the input data and applies an activation
% function to the weighted sum. This could be just a linear function or
% logistic function. In most cases, the activation function is a non-linear
% function, such as the rectified linear unit (ReLU) \notepw{glossary} or the
% sigmoid function. This is inspired from real neurons, that (in a simplified
% way) integrate the weighted information that they receive through their
% dendrites and then, when this weighted sum is high enough to cross a
% threshold, produce a spike. The activation function is equivalent to the
% threshold crossing.
%
% The weights of each neuron are adjusted using a process called
% backpropagation. Backpropagation is the algorithm that computes the gradient
% (the derivative) of the loss function with respect to the weights of the
% model. The gradient is then used to adjust the weights so that the loss
% function is minimized. The loss function is a function that measures the
% difference between the output of the model and the true output. The loss
% function is minimized by adjusting the weights of the model. If we stick to
% our example, during each step of the training loop, the model would consume
% an image. If the image is that of a dog and the models prediction is a
% probability of 20\% for the \enquote{dog} class and 60\% for the
% \enquote{cat} class, the loss function quantifies this error from the ground
% truth (which would be 100\% and 0\% respectively). Backpropagation then
% adjusts the weights of each neuron into a direction so that the model output
% comes closer to the ground truth. And this process is repeated many times.
%
% The simplest type of neural network is the multilayer perceptron (MLP)
% \notepw{glossary}. An MLP is a type of neural network that is composed of
% multiple layers of fully connected neurons. A fully connected neuron is a
% neuron that is connected to all neurons in the previous layer. For most real
% world tasks, simply adding fully connected hidden layers is not sufficient to
% learn a good model. The model will either not be able to learn the features
% that are important to predict the output or severely overfit to the training
% data. This is were different types of feature extraction layers come into
% play. For example in the case of time-dependent data such as climate data,
% time series data or audio data, networks with a capability to remember and to
% view current events in the context of previous events are needed
% (Long-short-term memory blocks or transformer blocks such as in GPT). In case
% of images, networks with a capability to extract spatial features are needed.
% Convolutional neural networks are a type of neural network that is
% specifically designed to process images. They achieve this by including
% multiple layers of 2D convolutions between the input and output layers. The
% convolutional layers are able to extract features from the input that are
% important to classify the input into the desired classes. This is similar to
% how the human visual system works as
% \textcite{hubelReceptiveFieldsSingle1959} discovered. They found that the
% visual cortex of cats contains neurons that are specifically tuned to detect
% edges, corners and other simple features. These neurons are then connected to
% other neurons that are specifically tuned to detect more complex features
% such as shapes and textures. Convolutional neural networks are able to learn
% these features from the input data and use them to classify the input into
% the desired classes.
%
%
% These feature extraction layers are meticulously designed (by humans or by
% other neural networks that learn to optimize the architecture of neural
% networks) to be able to learn to extract features specific to the kind of
% input data. This is the so-called \enquote{backbone} of the network. The
% backbone is the part of the network that is responsible for extracting the
% features from the input data. The backbone is then connected to a head that
% is responsible for the prediction of the output. Depending on the task, the
% head can be of different architectures. For example, in the case of a binary
% classification task, the head can be a single neuron that returns the
% probability of the input data to belong to one of the two classes. In the
% case of a multi-class classification task, the head can be a
% \enquote{softmax} layer that returns the probability of the input data to
% belong to one of the classes. In the case of a detection task, the head can
% be a set of neurons that return the coordinates of the bounding boxes of the
% detected objects and the probability of the input data to belong to one of
% the classes. In reality, the model architectures currently in use are much
% more complex than the fundamental building blocks introduced here with many
% different feature extration blocks and multipurpose heads that can compute
% the loss using multiple loss functions for multiple tasks. 
%
% With this, in mind, we can now proceed to the final approach that I took to
% use machine learning to detect chirps in recordings of multiple fish. The
% next section, we will introduce how convolutional neural networks can be used
% to detect chirps in recordings of multiple fish.

% section  Convolutional Neural Networks (end) \section{The Chirp Detection
% Pipeline} % (fold) \label{sec:the_chirp_detection_pipeline}
%
% In the following sections, we will dive into the chirp detection pipeline.
% The first model, the detector, is a convolutional neural network that locates
% chirps on spectrograms. It consists of a convolutional backbone and a head
% that not only predicts the class of the object but also the coordinates of
% the bounding box for each chirp. The second model, another much simpler
% convolutional neural network, is the classifier, which assigns the detected
% chirps to the emitting fish. The two models are embedded into a pipeline
% that, when run, automatically detects and assigns chirps in recordings of
% multiple fish. Broadly, the pipeline is divided into two stages: the
% detection stage and the assignment. Both stage use different neural network
% models to process the input data and each stage has its own feature
% extraction pipeline, inference step and post-processing step. The assignment
% stage builds on the output of the detection stage and provides the final
% output of the pipeline. The pipeline takes unstructured electrode grid
% recordings as input and returns a structured output that contains the
% detected chirps and the emitting fish as well as various parameters, such as
% chirp height and duration than can be used to further analyze the detected
% chirps.
%
% Before I we introduce the full pipeline, we will first discuss the datasets
% that I used, how I generated training data for each of the models, how I
% chose the architecture of the models and how I trained the models.
%
% section The Chirp Detection Pipeline (end) chapter Background (end)

\chapter{Methods} % (fold)
\label{chap:methods}

Following our discussion on the nature of chirps and their detection methods,
this chapter outlines the pipeline developed for chirp detection. This pipeline
comprises four key components: 

\begin{enumerate}
  \item \textbf{Feature extraction from grid recordings:} This 
        initial step involves generating spectrograms to identify chirps.
  \item \textbf{Chirp detection:} In this phase, chirps are detected on the
        spectrograms, undergoing various transformations and filters.
  \item \textbf{Feature extraction for chirp assignment}: This involves
        analyzing individual chirps on the spectrograms to extract features
        necessary for assigning them to specific fish.
  \item \textbf{Chirp assignment}: The final step reduces the assignment
        problem into a binary classification task, addressed with another
        neural network.
\end{enumerate}

The pipeline utilizes two separate feature extraction processes and two neural
networks to detect and assign chirps. Training and evaluating these networks
are important steps in the process. Therefore, before jumping into the
pipeline's specifics, this chapter will first detail the training process for
each neural network, including the data used for training.

\section{Chirp detection}
\label{sec:cpd}

The issue of chirp detection is where most of the complexity of the pipeline
lies and most of the training data had to be generated. The following sections
will explore the different strategies I used to identify chirps in a set of
simple recordings and how I used them to build a training dataset for complex
grid datasets.

\subsection{Datasets}
\label{sec:rawdata}

I used multiple datasets to extract training data for the object detection
model. The first dataset was recorded during competition experiments by
\textcite{raabElectrocommunicationSignalsIndicate2021} in \textit{A.
leptorhynchus}. They staged competitions of two fish over access to a shelter
in experimental tanks and recorded electric signals on an electrode grid as
well as spatio-temporal behaviors using an near-infrared camera. The recordings
were made with a grid of electrodes and fish were able to move freely in the
tank. This provided to be vital for generating a diverse training dataset:
While fish were competing, they produced a large number of chirps with a much
higher diversity in their characteristics as compared to stationary
chirp-chamber recordings. Additionally, the movement of the fish resulted in
much more challenging spectrograms. This illustrated early on how diverse the
chirps of \textit{A. leptorhynchus} can be and how movement increases the
complexity of chirp detection. The disadvantage of this setup is that if I
wanted to use this data to generate training data, computed spectrograms would
have needed to be annotated by hand to extract the chirps. And in an ideal
world, a dataset used to train an object detection model would contain multiple
thousands of images, each with multiple instances of the object of interest.
This is why I decided to add a second available dataset of chirp-chamber
recordings.\label{sec:ccdata}

The chirp-chamber recordings were produced from single individuals of
\textit{A. leptorhynchus} that were restrained to a PVC tube using mesh on both
ends, as described in \textcite{stocklEncodingSocialSignals2014}. The fish
were stimulated by a pair of electrodes that were placed on the left and right
side of the chamber. The two recording electrodes where positioned at the head
and tail of the fish. Because the stimulation electrodes are isopotential to
the field of the fish, the stimulus did not interfere with the EOD and hence
was not recorded on the recording electrodes. The fish were then stimulated by
a sine-wave mimicking the presence of a conspecific. The frequency of the sine-wave was varied between $\pm$ 30 \si{\hertz} around the EOD$f$ of the fish.
While this dataset included a large number of chirps, the diversity of chirps
was limited to mostly very short type 1 and type 2 chirps
\parencite{englerSpontaneousModulationsElectric2000a}.

Because of the lack of diversity in this dataset, I also included a third
dataset of various experiments that mostly included two fish interacting. The
dataset was recorded by \textcite{obotiWhyBrownGhost2023} and most fish were
spatially restricted to a certain degree, meaning that the strongest signal on
one of two recording electrodes was always the EOD of the fish of interest.
This dataset includes a large variety of different experiments, all of which
are documented in the original publication.

\subsection{Training data acquisition}
\label{sec:dettraindata}

The first step in training the object detection model was to generate a large
number of training images, i.e., spectrograms with chirps. As the datasets that
I had access to at the beginning of the project were very large, converting all
the data to spectrograms and annotating the spectrograms by hand was not
feasible. This is why I decided to start by modeling chirps in simulated grid
recordings. At least in theory, this provides the ability to generate an
infinite amount of training data without the need of manual annotation.

\subsubsection{Initial training on modeled chirps}

To train an object detection model on spectrograms, merely modeling chirps is
not enough. To get a realistic spectrogram, a simulated recording of multiple,
moving and communicating fish is necessary. In this section, I will describe
the process of generating such a simulated recording and how to extract
training data from it. The first step was to model the frequency excursions
that happen during chirps.

\emph{Simulating a chirp.} To model chirps, I first needed to understand and
parameterize the chirps that are present in the recordings. As the only dataset
that I had access to at the time was the competition dataset of
\textcite{raabElectrocommunicationSignalsIndicate2021}, I decided to use the
chirps that were present in this dataset to model chirps. As a first step, I
had to extract instantaneous frequency estimates of the EOD of a chirping fish.
To do this, I used the chirps that were already sparsely detected by the
previous chirp detection approaches as first estimates on where to look for
more. I only considered the chirps of fish with the highest baseline EOD$f$.
The chirps of the fish with the lower baseline EOD$f$ would cross the frequency
band of the upper fish, rendering them unusable. I used a set of heuristics to
extract the instantaneous frequency for each chirp, which are described in
detail in the \hyperref[appendix]{appendix}.

To model chirps, I then fitted a Gaussian frequency modulation to the centered
instantaneous frequency of each chirp. The Gaussian modulation had a height
$h$, a width $w$ and a kurtosis $k$. The standard deviation was then computed
from the width and kurtosis:

\begin{equation}
  g(t) = h \exp\left(-\frac{1}{2}\left(\frac{t - \mu}{\sigma}\right)^{2}\right)^{k} \quad \text{with} \quad
  \sigma = \frac{w / 2}{2 \log(10)^{k/2}}
\end{equation}

Where $t$ is the time, $\mu$ is the center of the chirp and $k$ is the
kurtosis. A kurtosis of 1 corresponds to a perfect Gaussian. In addition to
this rather simple model, I also experimented with mixtures of up to three
Gaussian modulations. To achieve realistic type 1 and type 2 chirps, at least
two Gaussian curves should be needed as these chirps often contain an
undershoot or down modulation of the frequency at the end. The third Gaussian
was introduced to allow facilitate degree of freedom. But as model complexity
increases, so do the fitting errors, which is why I decided to stick with the
simple Gaussian model. After fitting I then linearly interpolated the three
parameters of the monophasic model. I sorted all parameter groups by the width
of the Gaussian and then interpolated all parameters to fill the gaps in the
parameter space. This allowed modeling chirps of any duration, height and
kurtosis, limited by the minima and maxima of the real chirps in the dataset.

\emph{Simulating a chirping fish.}\label{sec:simulatingachirpingfish} Moving
from frequency space into the time domain, the next step involves simulating
the EOD of \textit{A. leptorhynchus}, containing the simulated chirps. The EOD,
a periodic signal with a fundamental frequency and harmonics, can be modeled as
a sum of sinusoids, if its Fourier spectrum is known. Utilizing the
\href{https://github.com/janscience/thunderfish/}{\texttt{thunderfish}}
package, we can easily model an EOD, including specific amplitudes and phases
for different taxa.

To simulate an EOD's evolution over time, I start by selecting a baseline
frequency between 600 and \SI{1100}{Hz} from a uniform distribution, reflecting
the typical range for \textit{A. leptorhynchus}. To introduce chirps,
parameters are randomly drawn from the previously described interpolated
Gaussian model and  the resulting frequency curve is added to the EOD$f$
vector. The simulation randomly determines the number of chirps, ensuring no
more than one chirp per second and a minimum interval of
\SI{200}{\milli\second} between chirps. To achieve this, a random number of
chirps (constrained by the maximum rate) is randomly positioned onto the full
frequency trace and chirps that are too close are discarded.  Notably, chirps
also involve a change in EOD amplitude, typically a slight reduction. This
effect is simulated by inverting the Gaussian used for frequency modulation,
adjusting its scale, and multiplying it with the EOD. The resulting amplitude
is randomly set between 100\% and 20\% of the EOD's maximum amplitude.
Reconstructing the EOD as a sum of sine waves using the EOD$f$ vector and
additional parameters for phases and amplitudes results in a semi-realistic
simulation of a chirping fish that oscillates between negative one and one.

At this point, the simulation resembles a recording one would get from a
chirp-chamber setup: A single fish that is not moving and producing chirps,
recorded on two electrodes positioned at its head and tail. This is not enough
to emulate the complexity of a grid recording of multiple, moving fish. So as a
next step, let's add movement to the artificial fish.

\emph{Simulating a moving fish.} An electrode grid reduces the spatial
information to a two-dimensional plane. Because the power of the EOD decays
with distance, it is informative about the distance of the fish to the
recording electrode. In such electrode grid recordings, the positions of fish
in two dimensional space are reflected by the distribution of the EOD
amplitudes on the electrodes. And when a fish moves, the amplitude of the EOD
on each electrode changes. To simulate this, I first simulated the
\enquote{true} movement of a fish over time, from which I then calculated the
distances to all electrodes and the corresponding \acrfull{AM}. As the exact
movement patterns do not matter for this simulation, I used a simple random
walk. The algorithm gets a starting position, a target duration and two
\acrfullpl{PDF}, one for the direction and one for the distance of a single
discrete step. Manipulating the probability densities as well as the number of
steps per time strongly influences how the simulated fish moves. The PDF of the
step distance was a Gamma distribution normalized to the number of steps per
second with a peak at \SI{0.2}{\meter\per\second}, which is equivalent to
velocities observed in the field. The direction PDF was a sum of two Gaussian
distributions with a peak at \SI{0}{\degree} and \SI{180}{\degree}. Increasing
the width of the Gaussians for the direction PDF would make the movement
trajectories much less smooth, as the fish would change direction more
frequently. And \emph{vice versa}, narrowing both Gaussians would make them
smoother. Increasing the weighting of the forward Gaussian, i.e., the one at
\SI{0}{\degree} relative to the backward Gaussian, would decrease the number of
times the fish moves backwards. The forward Gaussian was set to a standard
deviation of 0.2 and 0.1 for the backward Gaussian on an angular unit circle.
The PDF of the backward Gaussian was weighted as 20\% of that of the forward
Gaussian. The chosen parameters ought to simulate the characteristic
forward-backward movement of knife fish like \textit{A. leptorhynchus}. As a
result, we get a vector that describes the position of a fish over time in a
two-dimensional plane. Since the random walk often breached the boundaries of
simulated space, I included boundaries that the fish could not cross. The
resulting positions were \enquote{folded} back into the simulated space (figure
\ref{fig:space_folding}) after generating the random walk.

\begin{figure}
  \centering
  \includegraphics[width=\textwidth]{../figures/space_folding.pdf}
  \caption{Simulated fish movements for three fish. The circles indicate
  electrodes and the black border is the boundary of the virtual space. The
  positions are simulated by a random walker that is capable of stepping into
  any direction of a circle around it. The trajectory of each step underlies a
  bimodal PDF as real fish mostly move forwards or backwards. The
  distance of each step is drawn from a Gamma distribution limited by
  velocities estimated from fish in their natural habitat with a peak at
  \SI{0.2}{\meter/\second}. To maximize computational efficiency, the positions
  are first simulated and then folded back across the boundaries of the virtual
  space until they are contained in the boundaries.}
  \label{fig:space_folding}
  \vspace{0.5cm}
\end{figure}

From these $x$ and $y$ positions, the next step was to simulate the power of
the EOD on each electrode depending on the position of the fish. For this, I
defined a grid of 16 virtual electrodes spaced \SI{50}{\centi\meter} apart in
the simulated space (figure \ref{fig:space_folding}). In the previous steps, we
covered everything needed to simulate a fish that is recorded on a single
electrode. To simulate how this fish would be recorded on the 16 virtual
electrodes, we can simply duplicate the signal 16 times. We then
retrospectively adjust the AM so that it is proportional to the distance between the fish and each of the 16 electrodes. By using Pythagoras' theorem
and the known electrode positions as well as the known fish positions, we can
calculate, for each point in time, the distance between the virtual fish and
each virtual electrode. As the amplitude of the EOD decays with the squared distance, inverting, squaring and normalizing the distance gives us the
amplitude of the EOD$f$ on each electrode. Normalization is necessary, as the
simulated EOD is restricted between 1 and $-1$, while the distance
can become much greater. To obtain an amplitude modulation that we can add to
the EOD, this AM computed from the distances also has to be restricted to
between 1 and $-1$. The main assumption here is that the power of
the electric field decays equally in all directions around the fish. In
reality, the field of the fish is an oscillating dipole with two lobes. But
because I assumed, that this property is not reflected in the appearance of the
sum spectrogram across all recording channels, I ignored it for now. Adding
this distance transformed into an AM to every single EOD gives us a matrix that
contains the EOD including position-dependent AMs. The result is a simulated
electrode grid recording of a single, moving fish that produces chirps.

\emph{Simulating multiple fish.} A single fish is not enough to emulate the
complexity of a grid recording in the field. In reality, the frequencies of
fish often overlap or cross each other, not to speak of the chirps that are
produced. To simulate this, I simply repeated the process of simulating a
single fish a random number of times (between one and five fish
with a minimum $\Delta\text{EOD}f$ between two fish of \SI{20}{\hertz}).
Combining the separate fish is then simply the sum of the matrices that contain
the EOD over time on all electrodes. The result is a simulated grid recording
of multiple, moving fish that produce chirps.

\emph{Simulating noise.} The simulated recordings are very clean and do not
resemble the complexity of a real recording. I added noise at different stages
of the simulation pipeline. Before doing so, we first have to consider how this
translates to spectrograms as this is the only feature that the model is able
to consume. Adding noise to the EOD$f$ would make the frequency band of a
single fish on a spectrogram broader. Adding noise to the EOD however, would
change the \acrfull{SNR}, i.e., the contrast between the \enquote{background}
and the frequency band of a single fish on a spectrogram. These sources of
noise can be easily implemented, given that they are static over time. So as a
first step, I added band-limited Gaussian noise to the baseline EOD$f$ of each
fish (cutoff at 0.0 \si{\hertz} and 0.01 \si{\hertz}, $\sigma = 5
\si{\hertz}$). With the frequency trace of each chirp, I multiplied
band-limited Gaussian noise (cutoff at 600 \si{\hertz}, $\sigma = 1
\si{\hertz}$) as the noisiness of the instantaneous frequency appeared to
increase proportionally with the increase in frequency relative to the baseline
EOD$f$. To the waveform of an EOD I additionally added Gaussian noise ($\sigma
= 0.01$) to simulate its presence in the recordings of single fish. The part
where it gets complicated is the noise that is non-stationary, meaning that it
substantially changes over time. This kind of noise is particularly strong in
laboratory recordings as many appliances, like lights, cables, and other
equipment, introduce electromagnetic noise. Due to its heterogeneity, it is
impossible to simulate \textit{ad-hoc}. It depends on the building, the room,
the time of day and many other factors. So instead of simulating it, I used the
noise floor from real recordings and added it to the simulated recordings. To
achieve this, I selected areas from the Raab dataset which were known for
either contain little or no chirps. I scaled the simulated recording so that it
fits the units of the recorded snippet and added it to my simulated ensemble of
fish to get a realistically noisy backdrop. This still contains the EODs of the
real fish from the Raab dataset, but as they do not produce any chirps, they
can be considered as additional noise.

The result is an artificial grid recording that contains multiple, moving fish
producing chirps and that has a realistic noise floor. As I have simulated
them, I also know when which fish produced a chirp and what the chirp looked
like. But this is not yet a dataset that can be used to train an object
detection model. To train a detection model, we need spectrograms that contain
the frequency bands of chirping fish and labels that contain the coordinates of
bounding boxes around the chirps. So now is a good time to introduce the format
of the training dataset and how spectrograms are computed.

\emph{Spectrogram conversion.}\label{sec:spectrogramconversion} Computing
spectrograms is a computationally expensive task, especially with weeks of
recordings. The two factors that limit how much data of one recording should be
converted to spectrograms at once are the available memory and the time it
takes to compute. Using modern graphics cards can decrease computational time
but limits the amount of data in terms of memory, as \acrshort{GPU}s usually
have less memory than \acrshort{CPU}s. Using an NVIDIA GeForce RTX 4080 with 16
GB of memory, I was able to convert three minutes of recordings from 16
channels sampled at 20 \si{\kilo\hertz} at once (nfft=4096, hop length=410). Converting three minutes of recording takes
approximately 300 \si{\milli\second}. All channels in a grid recording are
converted to a sum spectrogram. Additionally, the spectrogram is transformed
from the power scale to the decibel scale. The windows that I extracted from
the recordings always had an overlap of one second. Another issue is that the
resulting spectrograms are quite large: Finding an event of a few tens of
milliseconds in three minutes of data is like finding a needle in a haystack.
Additionally, the relevant part of the frequency spectrum only reaches up to 2
\si{\kilo\hertz} of the available 10 \si{\kilo\hertz}. To make this data easier
to process for the neural network, I split the spectrograms into 15
\si{\second} chunks with 1 \si{\second} overlap. I also split them along the
frequency dimension into the ranges 0-1000 \si{\hertz}, 500-1500 \si{\hertz},
1000-2000 \si{\hertz}. To generate training data, each of these spectrogram
chunks is saved to disk as a greyscale PNG image.

\emph{Generating labels.} To train a model for predicting the position of
chirps on spectrograms, we need labels that allow the model to learn. A common
way to label the position of objects in images, or in this case, chirps on
spectrograms, is, to use bounding boxes. A bounding box is a rectangular box
around an object of interest that is defined by the coordinates of the four
corners. In the case of simulated data, these labels can be generated from the
known chirp times, widths and heights. To do this, I had to find a
transformation that approximates the width and height of a chirp on a
spectrogram given the same on the instantaneous frequency. As we have seen in
figure \ref{fig:multifish_chirps}, chirps look different on spectrograms
compared to their instantaneous frequency trace. I approximated the width of
the chirp on the spectrogram by adding a term to the original width that is
dependent on the sample rate $f_s$ and spectrogram: $\text{nfft} / f_{s} \times
0.8$. This is necessary because the minimum possible width of a chirp is the
width of the window that is used to compute the spectrogram. The upper boundary
of the bounding box was then set to the maximum frequency of the chirp and a
certain frequency padding given by $f_{s}/\text{nfft}\times5$. The lower
boundary was set to the minimum frequency of the chirp minus the same frequency
padding.

\emph{Generating the dataset.} The dataset is then generated by converting the
simulated grid recordings to spectrograms and saving them to disk as PNG images
and generating labels from the known chirp times, widths and heights. The
labels for each spectrogram were saved as \texttt{txt} files with the same
file name as the image file and coordinates transformed relative to the image
dimensions. This is the common dataset format for YOLO-style object detectors.
The result was the first dataset that I used to train the object detection
model. From now on, I will refer to this dataset as the \emph{simulated
training dataset}.

\subsubsection{Manual annotation with a human in the loop}

In addition to simulated data, real data was also used to train the object
detection model. To efficiently manage the annotation workload, subsets known
to contain chirps were extracted from the Raab datasets. These subsets were
pre-annotated using a faster R-CNN model, an established CNN for object
detection that was fine-tuned on simulated data
\parencite{renFasterRCNNRealTime2016}. The pre-annotations were imported into
\href{https://labelstud.io/}{\texttt{label-studio}}, a web-based, open-source
annotation tool designed for human-in-the-loop annotation processes
\parencite{tkachenkoLabelStudioData2022}. This tool facilitated the refinement
of pre-annotations and the addition of new ones by human annotators. The
refined annotations, from now on termed the \emph{recorded training dataset},
were exported and progressively expanded through cycles of model training,
pre-annotation, human correction, and export. The final dataset comprised 537
images of spectrograms with 14,832 annotations, averaging 27.6 chirps per
image.

\subsubsection{Hybrid simulations}

The last and probably most important dataset that I used to train the object
detection model is what I will refer to as the \emph{hybrid training dataset}.
This is a combination of a fully simulated grid recording with the difference
that the instantaneous frequency evolution of each chirp is not generated by a
simple Gaussian model as it was previously, but extracted from real recordings.
Previously, I introduced the \hyperref[sec:ccdata]{chirp-chamber dataset} which
has the advantage that the fish are solitary and stationary, facilitating the
extraction of the instantaneous frequency. Additionally, these recordings are
differential, meaning that they contain just a single channel signal. Chirps
were simply detected as peaks on the instantaneous frequency. I used a set of
hand-tuned heuristics to extract chirps from these chirp-chamber recordings
that are described in detail in the \hyperref[appendix]{appendix}. The number
of peaks discarded in each filtering step are visualized in figure
\ref{fig:chirp_baseline_smoothing} A.

This resulted in 9,082 extracted chirps. The resulting chirps still had an
often jittery baseline frequency. If I would put these chirps directly into the
simulation pipeline, the area of the spectrogram around the chirp would be
different from the rest of the baseline EOD$f$ band. This is a feature a neural
network could learn but that is not present in real data. To prevent this, I
had to find a way to make the baseline frequency, i.e., everything left and
right of the chirp, as smooth as possible while keeping the natural variability
of the EOD$f$ during a chirp.

\begin{figure}
  \centering
  \includegraphics[width=\textwidth]{../figures/plot_chirp_baseline_crossfade_smoothing.pdf}
  \caption{
    Visualization of the filter contributions and frequency post-processing
    steps. (A) Shows how many of the detected peaks in the EOD$f$ were
    discarded by each filter threshold. The thresholds were chosen to remove
    false positives and to keep the chirps that are most likely to be real. For
    details see \hyperref[appendix]{appendix}. (B) Illustrates how the baseline
    instantaneous frequency around a chirp was smoothed. I used the smoothed
    baseline to find the beginning and end of the chirp and applied a
    cross-fade to the instantaneous frequency to make the baseline around the
    chirp as smooth as possible while retaining the natural variability of the
    instantaneous frequency during a chirp. The result is a centered
    instantaneous frequency that can be used to simulate a chirping fish.
  }
  \label{fig:chirp_baseline_smoothing}
  \vspace{0.5cm}
\end{figure}

To achieve this, I used cross-faded concatenation of a smoothed baseline before
and after the chirp (figure \ref{fig:chirp_baseline_smoothing} B). First, I
estimated the baseline frequency EOD$f_{bl}$ by computing the median of all
values in the boundaries $b_{lower}=\min{\text{EOD}f(t)}$ and
$b_{upper}=\left|\min{\text{EOD}f(t)}\right|$. I then descended the left and
right slopes of the smoothed chirp until I either reached the baseline or the
sign of the derivative flipped, indicating that the slope goes up again. This
is not as straight forward for the end of the chirp, as a chirp often has an
undershoot. So here I descended the right slope until the sign of the
derivative flipped twice: The first time when the slope goes up again after the
undershoot and the second time when the jitter of the noisy baseline starts.
The result is a start and stop time of a chirp, which includes the undershoot
and. This can be used to extract the chirp from the centered instantaneous
frequency. To get a smoother baseline, I applied a Gaussian filter ($\sigma=60
\, \text{samples}$) to the instantaneous frequency. I then used the 120 samples
after the chirp start and the 120 samples before the chirp stop as the
cross-fade regions. In these regions, I combined the original instantaneous
frequency EOD$f(t)_{\text{orig}}$ with the smoothed instantaneous frequency
EOD$f(t)_{\text{smooth}}$ using a weighted mean with linearly increasing and
decreasing weights $w(t)$ for the cross-fade regions respectively:

\begin{equation}
  \text{EOD}f(t)_{\text{cf}} = \frac{1}{2} \left((1 - w(t)) \text{EOD}f(t)_{\text{orig}} + w(t) \text{EOD}f(t)_{\text{smooth}} \right)
\end{equation}

Finally, the baseline frequency was subtracted from EOD$f_{cf}$. The result
EOD$f_{\text{cf}}$ was now the instantaneous frequency of 9,011 chirps with a
perfectly smooth baseline that retained the natural frequency variability of a
real chirp. The remainder did not successfully pass the cross-fading and was
discarded. This happened when the descent of the slopes on either flank never
reached the baseline or a change of sign until the end of the snippet.

The resulting distribution of chirp heights and widths was severely skewed
towards small and short chirps as shown in figure
\ref{fig:chirp_overview_clustered_plot}. As I wanted the model to explicitly
learn to detect long and high chirps as well, I re-sampled the chirps to have a
uniform distribution by using a histogram equalization algorithm on both the
chirp heights and widths. This algorithm computed a 100-bin histogram for both
the distribution of chirp heights and widths. It then filled all bins by
duplicating random chirps from each bin to achieve a uniform distribution
across the bins. After histogram equalization, the dataset contained 121,380
chirps. Duplicating chirps is not ideal, as it introduces reduced variability
of large chirps. On the other hand, how a chirp appears on a spectrogram is
strongly dependent on the noise floor, how many other fish are present in the
recording and what their EOD$f$ is. So while the same chirp might appear many
times in the dataset, it will do so in many different contexts. This is much
better than having strong class imbalance. The resulting instantaneous
frequencies were then used as sources of chirps in the previously introduced
\hyperref[sec:simulatingachirpingfish]{simulation pipeline} to generate grid
recordings of multiple moving fish that produce chirps. Training data, i.e.,
spectrograms and labels, were then generated from these simulated grid
recordings as described in \hyperref[sec:spectrogramconversion]{spectrogram
conversion} and following sections. Figure \ref{fig:annotated_chirps} shows two
training samples. Each training sample is a single spectrogram image containing
many bounding boxes for chirps.

\begin{figure}
  \centering
  \includegraphics[width=\textwidth]{../figures/chirp_overview_clustered_plot.pdf}
  \caption{
    Extracted chirps for the hybrid training dataset. (A) Clustered
    instantaneous frequencies of the chirps for visualization. Chirps were
    clustered by fitting a \acrfull{GMM} by optimizing the \acrfull{BIC} on
    the first three principle components of the instantaneous frequencies. The
    upper panel shows the number of chirps per cluster. Below are the medians
    and 25\% and 75\% percentiles of the instantaneous frequencies of the
    chirps. Small and short chirps are disproportionally represented. (B) The
    distribution of chirp heights and widths as well as the troughs after the
    chirp indicated by the hue. The distribution is severely skewed towards
    small and short chirps. There is also no apparent correlation between chirp
    height and width and the trough after the chirp: Most small chirps have no
    trough, while large chirps, particularly those at around 300 \si{\hertz},
    have a trough. But chirps that are even larger or last longer have smaller
    troughs again.
  }
  \label{fig:chirp_overview_clustered_plot}
  \vspace{0.5cm}
\end{figure}

\begin{figure}
  \centering
  \includegraphics[width=\textwidth]{../figures/annotated_chirps.pdf}
  \caption{
    Two training samples from the hybrid grid datasets with real chirps. Each
    sample contains \SI{15}{\second} of the signal and captures a frequency
    range of \SI{1}{\kilo\hertz}. Bounding boxes are automatically computed by
    respecting the width and height of each underlying chirp as well as the
    \texttt{nfft} of the spectrogram to account for the artifacts introduced by
    estimating the spectrograms.
  }
  \label{fig:annotated_chirps}
  \vspace{0.5cm}
\end{figure}

\subsection{Model selection}

Deep learning and computer vision have seen a rapid development in the last
decade. This has led to a large number of different object detection models
that are tailored to different tasks and have different trade-offs in terms of
speed, accuracy and ease of use. The popular object detection models are the
R-CNN and YOLO families
\parencite{girshickRichFeatureHierarchies2014, redmonYouOnlyLook2016}. The
R-CNN family is based on the \emph{region-based convolutional neural network}
and consists of the original faster R-CNN, the R-FCN and the
Mask R-CNN. The YOLO family is based on the \emph{You Only Look Once}
model and consists of the original YOLOv1, the YOLOv2 up to
YOLOv9 and various deviations such as YOLONas and PP-YOLO,
which each have their own tweaks in model architecture and training process.

To start off, I decided to use the faster R-CNN model as it is well
documented and included in the \href{pytorch.org}{\texttt{PyTorch}} ecosystem.
The faster R-CNN model is a two-stage object detection model that first
generates a set of region proposals and then uses a region of interest pooling
layer to extract features from each region proposal. The features are then used
to predict the class and the bounding box of the object. I used the
ResNet-50 as the backbone. The features extracted by the backbone are
then used to generate the region proposals and to predict the class and the
bounding box of the object. In this case, I implemented the model training loop
using the \href{pytorch.org}{\texttt{PyTorch}} library. The model was trained
using the \texttt{Adam} optimizer with a learning rate of 0.001 and a batch
size of 8. This is the model used for the pre-annotation of the recorded
training dataset.

In addition to the faster R-CNN model, I also used the YOLOv8 and the
YOLOv9 models. These models are single-stage object detection models
that predict the class and the bounding box of the object directly from the
input image. For the YOLO models, implenting a training loop was not nessecary,
as the models were trained using the
\href{https://ultralytics.com}{\texttt{ultralytics}} library
\parencite{yolov8_ultralytics}. This includes many pre-trained models as well
as features such as a memory dependent dynamic batch sizes and a model
evaluation pipeline, which simplified the training process. All YOLOv8 as well
as the YOL0v9E models were trained and can be used as the object detection
backend in the chirp detection pipeline.

\subsection{Training \& hyperparameters}

In the final detection training dataset I combined the hybrid training dataset
with the human annotated dataset resulting in 2,159 training images. YOLOv8 and
YOLOv9 are the two main models that are used in the chirp detection pipeline,
as they are the most accurate, fast and easy to fine-tune. The models were
trained using the default hyperparameters of the \texttt{ultralytics} 8.1.23
library except for setting the image size to (800x800). Each model was trained
for 250 epochs and 5 cross-validation folds (figure
\ref{fig:yolo_crossval_training}). The batch size had to be decreased for the
larger models to fit into the memory of the GPU.

\begin{figure}
  \centering
  \includegraphics[width=\textwidth]{../figures/yolo_crossval_training.pdf}
  \caption{
    Training curves for all models that are trained for the detection step.
    Metrics were computed using 5-fold cross-validation for each model. The
    line corresponds to the median across the folds and the colored area to the
    \ordinalnum{25} and \ordinalnum{75} percentile respectively. In most case,
    the percentiles are barely visible despite training on a different subset
    of the dataset with the same random seed. Each row corresponds to a model
    and each column corresponds to a combination of loss metric and performance
    metric. The column header indicates which loss and which performance metric
    is plotted and the line colors indicate on which data (training /
    validation) the metric was evaluated. Losses decrease for all YOLOv8 models
    without too strong deviations between the training and validation loss.
    Losses and performance metrics of YOLOv9E also follow the predicted path
    but show strong fluctuations after 100 epochs are trained.
  }
  \label{fig:yolo_crossval_training}
  % \vspace{0.5cm}
\end{figure}

\subsection{Evaluation}
\label{sec:deteval}

To understand how object detection models can be evaluated, we first have to
define what a true positive and what a false positive is. In other words, how
close does a predicted bounding box have to be to the true bounding box to be
considered a true positive? And how far does a predicted bounding box have to
be from the true bounding box to be considered a false positive? To answer this
question, we can use the \acrfull{IoU} metric. The IoU between a ground truth
bounding box $y$ and a predicted bounding box $\hat{y}$ is defined as the area
of the intersection of the two bounding boxes divided by the area of the union
of the two bounding boxes:

\begin{equation}
  \text{IoU}(y, \hat{y}) = \frac{\text{area}(y \cap \hat{y})}{\text{area}(y \cup \hat{y})}
  \label{eq:iou}
\end{equation}

Traditionally, the IoU is used by thresholding it in two ways: The simple
approach is to assume that a predicted bounding box is a true positive if the
IoU is greater than a certain threshold, in most cases 0.5. A more nuanced
approach is to use a range of IoU thresholds and to compute the \acrfull{AP}
over the range of IoU thresholds, e.g., in between 0.5 and 0.95 in steps of
0.01. The average precision is defined as the area under the curve spanned by
the precision $p$ and recall $r$ as a function of the detection threshold
$\theta$, and is the most common metric to evaluate object detection models.
Precision and recall are computed from the number of true positives
($n_{\text{TP}}$), false positives ($n_{\text{FP}}$) and false negatives
($n_{\text{FN}}$):

\begin{equation}
  \begin{aligned}
    \text{p}(\theta) = \frac{n_{\text{TP}}(\theta)}{n_{\text{TP}}(\theta) + n_{\text{FP}}(\theta)} \quad \text{r}(\theta) = \frac{n_{\text{TP}}(\theta)}{n_{\text{TP}}(\theta) + n_{\text{FN}}(\theta)} 
  \end{aligned}
  \label{eq:ap}
\end{equation}

Naturally, the average precision at an IoU threshold of 0.5 (noted as
AP@.5) will result in a higher average precision than the average precision at
an IoU threshold of the previously mentioned steps (noted as
AP@[.5:.01:.95]). The latter can be formulated as the following:

\begin{equation}
  \begin{aligned}
    \text{AP@[.5:.01:.95]} = \frac{1}{N} \sum_{\theta}^{N} \text{AP}(\theta) \quad \text{where} \quad \theta \in [0.5, 0.95]
  \end{aligned}
  \label{eq:fancyap}
\end{equation}

\texttt{ultralytics} provides both metrics as the \acrfull{mAP}. This would be
the means of the average precision over all detection classes. As I am just
detecting a single class, the mAP is equivalent to the AP.

As the output of the model is a bounding box as well as the confidence for each
bounding box, I had to find the best confidence score to use as a threshold to
minimize false positives missed chirps. To do this, I used the F1 score, which
is the harmonic mean of the precision and recall.

\section{Chirp assignment}

In the assignment step, every detected bounding box needs to be attributed to a
fish in the recording. To assign each detected chirp to the identifier of the
emitter fish, I found the most informative feature to be the time dependent
power of the EOD$f$ of fish which intersect with the bounding box.

\subsection{Feature extraction}

To identify which frequency bands in chirps on spectrograms might contain the
EOD$f$ of fish, I estimated the baseline EOD$f$ of the emitter, which remains
relatively constant over short time spans. By summing the spectrogram along the
time axis to compute a power spectrum, I identified peaks representing the
EOD$f$ frequency bands of potential chirp emitters. Around each peak, a
\SI{10}{\hertz} wide band was used to calculate the average power within the
bounding box's duration for each frequency band. This feature is informative
because, during a chirp, the emitter's frequency increases, diverging from the
baseline and resulting in a power trough at the chirp's location due to the
frequency band shifting away from the baseline. However, this effect is only
partially observed due to the spectrogram's temporal resolution. While a higher
resolution would enhance this feature's utility, it would also reduce the
ability to distinguish between multiple fish due to overlapping frequency
bands. Despite this limitation, and the fact that this feature may be less
effective for detecting small chirps, it proved sufficiently reliable for
inclusion in the final pipeline. The feature extraction step is illustrated in
figure \ref{fig:assignment_feature_extraction} for a small and large chirp
respectively.

\begin{figure}
  \centering
  \includegraphics[width=\textwidth]{../figures/assignment_feature_extraction.png}
  \caption{
    Visualization of the feature processing steps starting from the predicted
    bounding boxes for a small chirp (A) and a big chirp (E) respectively. The
    power spectrum, i.e., a summation across the time axis of a bounding box,
    provides peaks at frequencies that could have been the fundamental
    frequencies of potential emitter fish (B, F). For each of these
    frequencies, the temporal evolution of the underlying power usually
    contains a trough if the current baseline frequency is that of an emitter
    fish (C, G). These power time series are then classified using a 1D
    \acrshort{CNN} (D, H). For each bounding box, the candidate frequency with
    the highest classification score is the fundamental frequency of the
    emitter fish.
  }
  \label{fig:assignment_feature_extraction}
  \vspace{0.5cm}
\end{figure}

\subsection{Training data acquisition}

The issue with this feature is that (1) deciding to which fish a chirp belongs
is not always clear on spectrograms and (2) the way this would be annotated is
not implemented in any annotation tools. So to train the assignment model, I
had to completely rely on my hybrid training dataset. To get the training data,
I computed the power of the fundamental frequency for each fish in each
bounding box during the extraction of the detection model training data. As the
chirps are simulated, the resulting dataset is truly a ground truth.

\subsection{Model selection}

To assign the chirps, I trained a simple binary classification model. For each
fish in a single bounding box, the model should return a score which indicates
the probability that a chirp belongs to the fish. Chirps are then assigned by
just accepting the fish with the highest score. As the model had to classify
simple one-dimensional arrays, I experimented with different types of
multi-layer perceptrons and 1D convolutional neural networks. The final model
is a simple CNN with two convolutional layers followed by four fully connected
layers, that are all followed by a \acrfull{ReLU} activation function except
for the output neuron, which is passed through a sigmoid. The model was build
and trained using the \href{pytorch.org}{\texttt{PyTorch}} package.

\subsection{Training \& hyperparameters}

The model was trained using the hybrid dataset. As I had much more
\enquote{non-emitters} than \enquote{emitters} in the dataset, I used a
weighted binary cross-entropy loss to balance the classes. The model was
trained with the \texttt{Adam} optimizer with a batch size of 32 and a learning
rate of 0.0001 over  11 epochs, after which the model converged (figure
\ref{fig:assignment_training} A).

\subsection{Evaluation}

To evaluate the model performance I used the \acrfull{AUC} of the \acrfull{ROC}
curve. The AUC is a metric that measures the ability of the model to
distinguish between the two classes. The AUC is calculated by plotting the
true positive rate (TPR) against the false positive rate (FPR) at various
threshold settings $\theta$ for the model confidence. The AUC is then the area
under this curve.

\begin{equation}
  \text{TPR}(\theta) = \frac{n_{\text{TP}}(\theta)}{n_{\text{TP}}(\theta) + n_{\text{FN}}(\theta)} \quad \text{FPR}(\theta) = \frac{n_{\text{FP}}(\theta)}{n_{\text{FP}}(\theta) + n_{\text{TN}}(\theta)} \label{eq:roc}
\label{eq:aucroc}
\end{equation}

The AUC ranges from 0 to 1, where 0.5 corresponds to a model that randomly
assigns the classes and 1 corresponds to a model that always assigns the
correct class. The AUC is a good metric to evaluate the performance of the
model as it is independent of the threshold setting and thus can also be used
to choose the best threshold setting for the model in production. As the AUC
can provide overly optimistic performance estimates for strongly imbalanced
classes, I also used the average precision, i.e., the area under the
precision-recall curve as well as the F1 scores (as explained in section
\ref{sec:deteval}) to evaluate the models performance and find the best
threshold.

\section{Behavioral analyses}

All behavioral analyses on the Raab dataset were conducted as described in
\textcite{raabElectrocommunicationSignalsIndicate2021}. The agonistic
interactions were annotated using
\href{https://github.com/olivierfriard/BORIS}{\texttt{BORIS}}, a free and
open-source software designed to facilitate the coding of behaviors in animal
behavior studies. The EOD$f$s of the fish were tracked using the
\href{https://github.com/tillraab/wavetracker}{\texttt{wavetracker}} package
for \texttt{Python}. Annotated behaviors and tracked EOD$f$s were already
available. I detected and assigned chirps to individual fish using the here
introduced
\href{https://github.com/weygoldt/chirpdetector/}{\texttt{chirpdetector}}
package. The chirp rates during the time of agonistic events were computed by
convolving the chirps of each fish with a Gaussian kernel with a standard
deviation of one second and a height of one. The \ordinalnum{1} and
\ordinalnum{99} confidence intervals were estimated by jackknife resampling
corresponding of an $\alpha$-level of 1\%. The null hypothesis, which is that
the chirp rates are independent of the on- and offset of agonistic events was
tested using bootstrapping. These methods are introduced in detail in
\textcite{raabElectrocommunicationSignalsIndicate2021}. Statistical analyses of
chirp counts and chasing durations were performed using \acrfullpl{GLMM}
implemented in the
\href{https://cran.r-project.org/web/packages/lme4/index.html}{\texttt{lme4}}
package in \texttt{R v4.3.3} with the pairing ID as a random factor
\parencite{batesFittingLinearMixedEffects2015}. To test whether the competition
status of two fish had an effect on the number of chirps each fish produced, I
used a GLMM of the Poisson family with a logarithmic link function. To model
the effect the duration of chasing had on whether loser produced chirps or not,
I weighted each data point (duration of a single chasing and the presence of
chirps) by (1) how many chasing events were produced by each pairing and (2)
the fraction of chasings in which chirps were produced, to not bias the model
to a specific trial. I then used a Binomial GLMM with a Logit link function.
The effect size is reported as Cohen's $d$. 


\chapter{Results} % (fold)
\label{chap:results}

The role of chirps as communication signals of wave-type weakly electric fish
is still mysterious, partly because they are not easy to study. I have now
introduced two neural networks that can be used to detect and assign chirps to
their emitter. In the following chapter, we see how these models can be
evaluated and integrated into a pipeline to automate chirp detection. This will
allow us to detect chirps on a large dataset from a behavioral experiment to
explore relationships between chirps, dominance status, and agonistic
interactions in competitive contexts of \textit{A.~leptorhynchus}.

\section{Architecture of the chirp detection pipeline} % (fold)
\label{sec:the_pipeline}

The chirp detection pipeline includes two feature extraction and two inference
phases. Initially, spectrograms are extracted from EOD recordings and are
passed into the detection model that returns time and frequency ranges, where
it detected chirps. Following chirp detection, a second feature extraction
phase gathers the temporal evolution of the powers of potential emitter
EOD$f$s, which a subsequent model classifies in the final inference step. This
way, each detected chirp is assigned to a fish. The full pipeline can be
accessed as a \texttt{Python} package and all code necessary to detect chirps
is freely available on GitHub
(\url{https://github.com/weygoldt/chirpdetector}). The upcoming sections detail
all processing steps and the models involved.

\subsection{Detection feature extraction}

The recordings from electrode grids have as many channels as there are
electrodes. The initial feature extraction phase transforms raw electric
recordings into spectrograms using a sliding window approach on every channel.
These spectrograms are then combined into a sum spectrogram and converted to
the decibel scale. Optimized for handling large datasets, these spectrograms
are computed \enquote{on-the-fly} and discarded post-detection model analysis.
To facilitate detection of small objects by the model and manage computational
costs effectively, spectrograms are segmented into smaller tiles
\parencite{huangTilingStitchingSegmentation2019}. Each tile spans
\SI{15}{\second} with \SI{1}{\second} overlap, and covers a frequency range of
\SI{1}{\kilo\hertz} with a \SI{500}{\hertz} overlap up to a maximum frequency
of \SI{2}{\kilo\hertz}. For instance, a minute-long recording yields four
time-axis windows and three frequency-dimension tiles, totaling 12 spectrograms
processed concurrently.

These spectrograms are then analyzed by the detection model to identify chirp
bounding boxes, with parallel processing capacity constrained only by GPU
memory. On a NVIDIA GeForce RTX 4080 with 16 GB of graphical memory, the system
can process three minutes of data from 16 electrodes sampled at
\SI{20}{\kilo\hertz} simultaneously. All performance-relevant parameters are
adjustable depending on the available hardware.

\subsection{Chirp detection}

Inference is performed using a YOLOv8 model, encapsulated by a wrapper to
standardize output formats, such as bounding boxes and confidence scores,
across different model architectures. This wrapper, based on an abstract base
class, facilitates easy model interchangeability and allows for the
incorporation of new models as they become available.

I tested the performance of all YOLOv8 versions as well as YOLOv9E. Both
training and cross-validation indicated that YOLOv8XL outperforms others in
precision, recall, and F1 score (figure \ref{fig:detector_performance} A-C).
Precision, informative about how many of the  predicted chirps are actually
real chirps, (equation \ref{eq:ap}), reaches above 99\% for the YOLOv8XL
(figure \ref{fig:detector_performance} A). Recall, relating to how many of the
chirps that are truly present in the dataset, are also detected (equation
\ref{eq:ap}), also surpasses 99\% for YOLOv8XL (figure
\ref{fig:detector_performance} B). Together, this shows that most of our
predictions are trustworthy and that most chirps in the dataset will be
detected. Only approximately 1\% of detections will be false positives and only
approximately 1\% of existing chirps will be missed. Optimal detection
thresholds to balance precision and recall are identified at around 50\% for
all models by the F1 score (figure \ref{fig:detector_performance} C). YOLOv8XL
was leading in average precision that reached up to 86\% (figure
\ref{fig:detector_performance} D, E, equation \ref{eq:ap}). The average
precision is a measure combining precision and recall into a single metric to
make comparisons of models easier. This is lower to the previous metrics as
here, the average precision is evaluated at AP@[.5:.01:.95] while the curves
are evaluated at a single IoU threshold of 0.5 (see section \ref{sec:deteval}).
Still, AP@[.5:.01:.95] of over 85\% can be regarded as a very precise object
detector, surpassing the current state of the art benchmarks for multi-class
detectors \parencite{wangYOLOv9LearningWhat2024}. Surprisingly, YOLOv9E, shows
subpar performance despite beeing the newest model, potentially due to training
instabilities (figure \ref{fig:yolo_crossval_training}). The performance
metrics of YOLOv8M to XL are closely matched, with average precisions exceeding
85\%. These medium sized models may suffice for less powerful devices or
embedded systems.

\begin{figure}[ht!]
  \begin{center}
    \includegraphics[width=1\textwidth]{../figures/kfold_modelsweep_performances.pdf}
  \end{center}
  \caption{
    Performance metrics for detection models on the validation dataset,
    including precision, recall, and F1 score. The F1 score serves as a
    performance indicator to determine the optimal confidence thresholds.
    Results, shown as median values from 5-fold cross-validation with
    \ordinalnum{25} and \ordinalnum{75} percentiles, though variations are
    minimal. Plots focus on areas where model performances diverge. Model
    precision quickly stabilizes near 100\%, with YOLOv8N slower to plateau
    (A). Model recall remains steady up to a 0.8 confidence, then declines (B).
    YOLOv8M-XL exhibit higher recall, and are less affected by the confidence,
    with YOLOv8XL leading. Optimal F1 scores occur around a confidence level of
    0.5 for all models (C). Larger models score higher, with YOLOv8XL achieving
    the top F1 score. Precision-recall curves, based on an IoU threshold of
    0.5, largely overlap, with YOLOv8XL showing the greatest area under the
    curve (D). Average precision, calculated across IoU thresholds from 0.5 to
    0.95, indicates overall model efficacy (E). Larger models outperform, with
    YOLOv8XL leading in average precision. YOLOv9E performs comparatively
    poorly and exhibits strong variability, probably due to the instabilities
    during training illustrated in figure \ref{fig:yolo_crossval_training}.
  }
  \label{fig:detector_performance}
  \vspace{0.5cm}
\end{figure}

Detected bounding boxes undergo several post-processing steps: Initially,
low-confidence detections are filtered out through thresholding. Subsequently,
non-maximum suppression eliminates overlapping boxes using the \acrshort{IoU}
metric, effectively reducing false positives and ensuring each chirp is
represented by a single bounding box. These bounding boxes mark detected
chirps, but we still do not know, which of the fish emitted these chirps. This
is why we need a second processing stage that assigns the detected chirps to
the emitting fish.

\subsection{Feature extraction for chirp assignment}

The second feature extraction, applied to each detected bounding box, begins by
estimating the fundamental frequency of all the fish potentially emitting the
chirp. This involves first interpolating the spectrogram within the bounding
box to $100 \times 100$ samples and calculating the power spectral density by
summing the resulting matrix along the time axis. Peaks in this spectrum
indicate the fundamental frequencies of potential emitters. A \SI{10}{\hertz}
window is used around each peak to sum power values in the spectrogram across
the frequency axis. The result is array representing the power of the
fundamental frequency over time for each potential emitter within a bounding
box. This process identifies deviations in the power of the fundamental
frequency associated with chirp emissions, where the power at the baseline
frequency drops. In non-chirping instances the power remains approximately
constant. The whole process is visualized in figure
\ref{fig:assignment_feature_extraction}. This feature extraction is efficient,
leveraging already computed spectrograms rather than having to go back to raw
electric recordings.

\begin{figure}[t]
  \begin{center}
    \includegraphics[width=\textwidth]{../figures/assignment_training.pdf}
  \end{center}
  \caption{
    Training curves and performance evaluation of the assignment model across
    5-fold cross-validation. The lines correspond to the median across the
    cross-validation folds and the colored areas to the \ordinalnum{25} and
    \ordinalnum{75} percentile respectively. The training as well as the
    validation loss converges after 11 epochs (A). F1-scores peak at a model
    confidence of approximately 82\% (B). The average precision of the model,
    i.e., the area under the precision-recall curve reaches approximately 97\%
    (C). The AUC of the ROC curve (D) reaches 99\% but a ROC evaluation has
    been shown to be overly optimistic when datasets are imbalanced
    \parencite{brancoSurveyPredictiveModelling2015}.
  }
  \label{fig:assignment_training}
  \vspace{0.5cm}
\end{figure}

\subsection{Chirp assignment}

For each bounding box, features are passed to the assignment model, which
evaluates the probability of each power array representing a chirp emitted by a
fish. The array with the highest probability identifies the emitter, assigning
its frequency as the fundamental frequency at that moment. The assignment
process, encapsulated by a wrapper, facilitates feature extraction and model
inference, with an abstract base class defining the model interface. Similar to
the detection model, this setup allows for easy model interchangeability and
simplifies the integration of new models into the pipeline.

The assignment model demonstrates high accuracy, achieving an F1-score of 91\%,
average precision of 97\%, and a AUC of 99\%, indicating effective chirp
assignment in most instances (figure \ref{fig:assignment_training} B-D). To
better understand the model's performance, I visualized the decision boundary
by reducing the 100-dimensional power traces to two dimensions using a simple
auto-encoder model (figure \ref{fig:assignment_zspace}). This reduction was
compared to results from \acrshort{PCA}, showing similar outcomes. The
visualization reveals a transitional zone between arrays classified as chirp
emitters and non-emitters, where ambiguous cases, likely representing short
chirps with less noticeable power drops, are found. Chirps with distinct power
reductions are clearly separated. The decision boundary, depicted with colors
for intermediate prediction scores, resembles a parabola rotated to
\SI{45}{\degree}. Most cases fall exactly onto the decision boundaries as
illustrated by the marginal distributions. This is surprising, given the high
performance metrics of the model. It indicates that the two-dimensional
representation is likely too simple to represent the true multidimensional
decision boundary and supports the used of a neural network, as a conventional
machine learning classifier running merely on the principle components probably
would have failed. But this remains to be verified.

\begin{figure}[!t]
  \begin{center}
    \includegraphics[width=\textwidth]{../figures/z_space_assignment_nn.pdf}
  \end{center}
  \caption{
    A visual simplification of the CNN classifier's decision boundary to a
    two-dimensional view for the evaluation dataset. Each point represents a 2D
    latent space mapping of power over time for a potential chirp emitter
    frequency, derived using a fully connected auto-encoder neural network
    trained on 20\% of the dataset. Points are colored based on the binary
    classifier's score, revealing a clear emitter class cluster that
    transitions into the non-emitter class. Insets highlight specific samples
    and their 2D mappings, illustrating a clear separation with definite
    emitters (top right) and non-emitters (bottom left) based on peak presence.
    Marginal distributions on each axis correspond to the chirp emitters and
    non-emitters after thresholding, based on an ideal 82\% score threshold
    from the F1-curve (figure \ref{fig:assignment_training}). These show
    overlapping distributions at the decision boundary. Despite this overlap,
    the classifier achieves an impressive 97\% average precision, as detailed
    in figure \ref{fig:assignment_training}.
  }
  \label{fig:assignment_zspace}
  \vspace{0.5cm}
\end{figure}

\newpage
\subsection{Postprocessing}

The two-step inference pipeline returns bounding box coordinates and predicted
fundamental frequencies of chirp-emitting fish. These outputs are integrated
with the \texttt{wavetracker} package, which tracks fish fundamental
frequencies over time on spectrograms and assigns identities to each frequency
track \parencite{raabAdvancesNoninvasiveTracking2022}. This allows
for estimating the number of recorded fish, approximating their positions, and
maintaining their identities. By matching the detected chirps' fundamental
frequencies with the closest fish frequency at the chirp time, chirps are
attributed to specific fish. Detections not aligning with a frequency track,
such as chirps detected on harmonics, are discarded.

The refined dataset contains chirp times, their dimensions on the spectrogram,
and the emitting fish's identity. Analysis of this dataset from staged
competitions revealed notable temporal patterns between chirps and aggressive
interactions, discussed in the following section.

\section{Evaluation of the detected chirps}

The chirp detector was applied to the Raab dataset, as published in
\textcite{raabElectrocommunicationSignalsIndicate2021}, along with additional
post-publication competition trials. This encompassed 69 trials in total, where
two fish competed for shelter over six hours in each trial, including three
hours of their active phase. To enhance signal clarity, a conservative
detection threshold of 0.7 was chosen, likely increasing the number of missed
chirps but minimizing false positives (figure \ref{fig:detector_performance}).
Across a total of 17.25 days of competition data, the detector identified
64,050 chirp instances.

Before exploring the potential role of chirps in the competitions documented by
\textcite{raabElectrocommunicationSignalsIndicate2021}, let us first assess
chirp detection effectiveness in new data through some examples. The detection
system successfully identifies and assigns most chirps, even when they
overlap other baseline EOD$f$ frequencies (figures \ref{fig:chirp_detections_simultaneous}
and \ref{fig:chirp_detections_bigchirps}). However, challenges arise in
specific scenarios. For instance, figure
\ref{fig:chirp_detections_simultaneous} illustrates a scenario where rapid
chirp production by one fish leads to detection failures and unreliable
assignments due to the spectrogram's limited temporal resolution. While such
high chirp rates are uncommon, they could hold significant behavioral interest.
Additionally, detecting chirps produced simultaneously by two fish poses
difficulties when their frequency difference is small, which causes chirps to
merge on the spectrogram. This often results in neural networks misidentifying
them as a single, large chirp. Despite the initial assumption that those cases
must be rare, figure \ref{fig:chirp_detections_simultaneous} presents two
instances of simultaneous chirps; one correctly assigned, and the other
falsely detected as a single chirp.

\begin{figure}[t]
  \begin{center}
    \includegraphics[width=\textwidth]{../figures/chirp_detections_simultaneous.pdf}
  \end{center}
  \caption{
    Detected and assigned chirps on an example spectrogram. Black lines
    correspond to the tracked fundamental frequencies of two fish. White boxes
    surround all detected chirps. The predicted fundamental frequency of the
    chirp emitter is illustrated by the dot in each bounding box. Overall, most
    chirps are detected well. The pipeline fails when chirps happen in too
    rapid successions to be separable on the spectrogram or when fish chirp
    simultaneously. The spectrogram was estimated using the same parameters
    used for chirp detection (nfft=4096, hop length=410) and interpolated
    (bicubic).
  }
  \label{fig:chirp_detections_simultaneous}
  \vspace{0.5cm}
  % 
\end{figure}

\begin{figure}[!h]
  \begin{center}
    \includegraphics[width=\textwidth]{../figures/chirp_detections_bigchirps.pdf}
  \end{center}
  \caption{
    Assignment can completely fail for large and long chirps and produces false
    assignment for large and short chirps. In one case, a large chirp is not
    detected at all. But the majority is detected and assigned correctly. The
    spectrogram was computed as described in figure
    \ref{fig:chirp_detections_simultaneous}.
  }
  \label{fig:chirp_detections_bigchirps}
  \vspace{0.5cm}
  % 
\end{figure}

Detecting large chirps presents another challenge. Although the training
dataset emphasizes these through sample duplication in histogram equalization,
long chirps are still accurately detected (figure
\ref{fig:chirp_detections_bigchirps}). However, assigning these chirps to the
correct sender often fails. This issue may arise because long chirps gradually
return to the baseline EOD$f$, which falls outside the detected bounding boxes.
Since emitter frequencies are determined from the power spectrum within these
boxes, long chirps might not exhibit distinct peaks at the emitter's EOD$f$.
Essentially, the power dip at the baseline is so broad that the baseline
frequency component becomes too faint to identify.

Another issue remains assignment in general: While most large chirps are
correctly assigned despite strongly overlapping with the baseline EOD$f$ of
another fish in figure \ref{fig:chirp_detections_bigchirps}, it includes three
examples of false assignments. Reasons could be amplitude troughs that
coincidentally occur due to noise or movements of the bottom fish, and that
resemble the troughs that happen during chirps. To conclude, while the great
majority of chirps are detected well, small issues remain. In most cases
however, they appear to be fixable.

Now let us take a look at all the bounding boxes that were predicted on the
dataset and are visualized in figure \ref{fig:all_boxes}. Predominantly, these
bounding boxes encapsulate small chirps, with the majority measuring below
\SI{200}{\hertz} in height and spanning \SI{100}{\milli\second} to
\SI{200}{\milli\second} in width (figure \ref{fig:all_boxes} B1). Remember that
the true chirp height and width does not directly map well to the dimensions on
the spectrogram: The most commonly observed small chirp is usually just
\SI{20}{\milli\second} in width but will be much wider on the spectrogram,
dependent on the parameters the spectrogram was computed with. This is the
reason of the compression of the width axis and the lack of chirps below
\SI{100}{\milli\second} in width.

Consistent with the training data in figure
\ref{fig:chirp_overview_clustered_plot}, the vast majority of chirps fall below
\SI{200}{\hertz} in height, as shown by the marginal distribution in figure
\ref{fig:all_boxes}. Another pattern that we have already observed in the
training data is that larger chirps tend to have stronger post-chirp
undershoots (figure \ref{fig:all_boxes} B2). But this pattern only extends to
short chirps, as the few long chirps that were detected do not show such strong
troughs (figure \ref{fig:all_boxes} B4).

In addition, two areas are visually separated from the main point cloud: Area
B3 is shaped similarly to area B2 but contains much less data points and they
appear to be moved up by \SI{100}{\hertz}. The simplest explanation would be
that these are large chirps produced by the fish with the higher EOD$f$, which
got falsely assigned because they overlapped with the frequency of the fish
with the lower EOD$f$. This way, the chirp height is incremented by the
frequency difference between the two fish, which often lies in a range between
\SI{50}{\hertz} and \SI{100}{\hertz}. Three examples for such false assignments
of large chirps are indicated in figure \ref{fig:chirp_detections_bigchirps}.
The second indicator supporting this assumption is the fact that this point
cloud has weaker post-chirp undershoots compared to area B2 in figure
\ref{fig:all_boxes}. This would also be explainable by the false assignment:
The trough is computed by subtracting the predicted baseline by the lower edge
of the bounding box. If a fish with a lower EOD$f$ is falsely assumed to be the
emitter, this difference becomes smaller.

\begin{figure}[!t]
  \begin{center}
    \includegraphics[width=\textwidth]{../figures/plot_chirp_bbox_distributions.pdf}
  \end{center}
  \caption{
    Statistics of all detected bounding boxes on the Raab dataset. Bounding
    boxes are centered at the baseline EOD$f$ of the fish and colored by
    density of the box type (A). Densities are estimated by Gaussian kernel
    density estimate on the first principle component of the height, width and
    trough parameters. The great majority of chirps are surrounded by bounding
    boxes with a width of \SI{100}{\milli\second} to \SI{200}{\milli\second}
    and a height below \SI{200}{\hertz}, which corresponds to small or type 2 
    chirps (B1). Post-chirp undershoots of the frequency become increasingly
    pronounced particularly for large but narrow chirps (B2). Chirps with a
    width beyond \SI{250}{\milli\second} are sparsely represented (B4). Beyond
    a height of \SI{500}{\hertz}, points appear to be a part of the point cloud
    below but incremented along the height dimension by approximately
    \SI{100}{\hertz}, which could be due to erroneous EOD$f$ predictions (B3).
    Chirps at around \SI{200}{\hertz} in height include a surprisingly narrow
    horizontal band that extends to durations of up to almost
    \SI{400}{\milli\second} (B5). These are most likely bounding boxes around
    chirps that happened in rapid succession, so that reliable chirp separation
    on the here used spectrogram parameters was not possible anymore. Chirps
    at a height range corresponding to this band are the only ones observed to
    occur in such rapid succession in this dataset, which would support this
    assumption.
  }
  \label{fig:all_boxes}
  \vspace{0.5cm}
\end{figure}

The second visually distinct area corresponds bounding boxes that are at
approximately \SI{200}{\hertz} in height but form a distinct band that extends
along the width axis (figure \ref{fig:all_boxes} B5). These are most likely
chirps that happen in such a rapid succession that they are inseparable on the
spectrogram. Given the width of \SI{100}{\milli\second} of most chirps in
figure \ref{fig:all_boxes}, rapid chirps that have a smaller inter-chirp
interval this, cannot be separated on these spectrograms. The frequency range
in which they occur supports this assumption, as chirps that are larger than
\SI{200}{\hertz} have not been observed to happen in such rapid succession
here. In some cases, the network still detects them as chirps but detection
boundaries can become much wider as they usually encompass multiple chirps at
once. Another interesting feature of this area is that it contains many chirps
with strong undershoots. This is especially true for more realistic (smaller)
chirp widths in this range. The most likely explanation is, that these are big
chirps produced by fish with a lower EOD$f$, which got falsely assigned to the
fish with the higher EOD$f$. This would drastically increase the trough while
decreasing the chirp height. Usually, the post-chirp trough is not stronger
than $1/3$ of the chirp height. In this area, the troughs are larger than the
chirp heights, which is not realistic. A simple filter that limits the
\enquote{chirp height to trough ratio} could remove these approximately 8,000
suspicious data points. Despite these two areas, most of the 73,134 points fall
into areas that are expected. In combination with the performance metrics, the
two example spectrograms, and hundreds of additional hand-checked spectrograms
the pipeline seems to detect chirps well enough to use the detections to
conduct behavioral analyses.

Now that we have demonstrated that the pipeline not only produces promising
performance metrics but that these metrics are also reflected on real,
previously unseen data, we can move on to analyze which fish produced chirps
and when they produced them.

% \newpage
\section{Subordinates chirp the most}

\textcite{raabElectrocommunicationSignalsIndicate2021} staged competitions
between two fish for access to a shelter. During these competitions, fish
partook in ritualized fights and chasing. In addition to video recordings, they
also recorded the electric fields of both fish using electrode grids. As video
recordings are not available for a large part of the dataset, and not annotated
yet for another part of the dataset, only 19 unique pairings between two fish
have annotated competitive behaviors. Figure
\ref{fig:competition_trial_overview} shows an overview of the available data
across one of these competition trials. A single trial consisted of three hours
of darkness, where fish were active and usually competed. During the subsequent
three illuminated hours, fish usually became inactive and sought shelter. The
fish that stayed in the superior shelter at the onset of the inactive period
was decided to be the winner of the competition or the \enquote{dominant}. In
the example trial, particularly the competition loser produced chirps, all of
which appear to visibly correlate with the chasing events.

To quantify whether this observation is systematically observable, I compared
the total number of chirps produced by dominants and subordinates. Overall,
winners chirped much less compared to competition losers (p<0.0001, d=66.5), a pattern that could be observed for all trials (figure \ref{fig:stats} A).
Losers produced chirps in an average of 29\% $\pm$ 23\% of all chasing
events.

As suggested in the literature, chirps could have been used as a submissive
signal by competition losers to stop extended chasing by the competition winner
or decrease the probability of future attacks by winners
\parencite{henningerStatisticsNaturalCommunication2018}. Alternatively, they
could serve as a signal of stress that is simply caused by extended chasing
without influencing the behavior of the winner. If chirps communicated
submission, leading to offset of chasing, the prediction would be that chasing
events were chirps were produced should have been shorter. This is not the
case, as the opposite could be observed: Chasing events, in which losers
produced chirps lasted significantly longer (p<0.0001, d=19.3). In every trial
in which the loser produced chirps during chasing, chirps were produced when
the median of the chasing durations was increased (figure \ref{fig:stats} B). 
An influence of chirping on the latency between the current and the next
chasing could not be observed. This suggests that competition losers started to
produce chirps as a response to beeing continuously chased by the competition
winner.

To understand with which particular behavior the chirps produced by the
competition losers or winners correlated temporally, I approximated how the
chirp rate changes during these events. As already implied in figure
\ref{fig:competition_trial_overview}, chirping of the competition loser but not
the competition winner was closely tied to chasing events (figure
\ref{fig:stats} C \& D vs E \& F). Chirp rates of the competition loser
significantly increased at the onset of being chased by the competition winner
(figure \ref{fig:stats} E). This effect is even more pronounced during the time
at which the chasing seizes: Here, the chirp rate of the competition loser
almost tripled and remained elevated for approximately \SI{10}{\second} beyond
the offset (figure \ref{fig:stats} F). This effect could not be observed for
competition winners, which produced little chirps and if so, showed no
correlations with the onsets or offsets of chasing events (figure
\ref{fig:stats} C, D).

\begin{figure}
  \begin{center}
    \includegraphics[width=\textwidth]{../figures/competition_trial_overview.pdf}
  \end{center}
  \caption{
    Overview of a single competition trial. The spectrogram illustrates the
    evolution of the discharge frequencies of both fish. Fish were left to
    compete for three hours while the light was off. This is also the time in
    which most of the frequency modulations are visible on the spectrogram. As
    the light is turned off, frequency changes of either fish become much
    rarer. Event plots show the times of the two annotated spatio-temporal
    interactions between the fish, the chasing and contact events. The majority
    of agonistic events are chasing events. Below are the event plots for the
    communication signals produced by either fish, separated by the competition
    status of the respective fish. Most rises but also most chirps are produced
    by the fish that lost the competition and did not get access to the
    shelter. Additionally, bouts of subsequent chirps produced by the loser
    visibly correlate with the chasing events.
  }
  \label{fig:competition_trial_overview}
  \vspace{0.5cm}
\end{figure}

\begin{figure}
  \begin{center}
    \includegraphics[width=\textwidth]{../figures/plot_behavioral_correlations_of_chirps.pdf}
  \end{center}
  \caption{
    Statistical summary of the role of loser chirps during ritualized chasing
    events. Points are connected according to the unique fish pairing in which
    they were observed. Competition losers produced more chirps than the
    competition winners (A). If the loser produced a chirp, chasing events
    lasted longer (B). Most chirps of the loser are produced close to
    the offset of the chasing event (figure \ref{fig:stats} F), which suggests
    that it is more likely that being chased for a longer time causes the
    losers to chirp. Chirp rates of competition losers centered around the
    onset and the offset of chasing events respectively (C-F). The chirp rate
    was estimated by convolving a Gaussian kernel with a binary chirp train.
    The confidence intervals are estimated by jackknifing and the baseline by
    bootstrap resampling. Chirp rates of the competition losers significantly
    increase after a chasing events starts (E). Specifically, most of these
    chirps that increase the chirp rate are precisely produced around the
    offset of chasing events (F), where the chirp rate almost triples.
  }
  \label{fig:stats}
  \vspace{0.5cm}
\end{figure}

% chapter Results (end)

\chapter{Discussion}
\label{chap:discussion}

The results indicate that chirps might play a role in mediating competitive
interactions among Brown Ghost Knifefish. By developing and implementing a deep
learning pipeline for automated chirp detection, I was able to observe and
analyze over 73,000 chirps during unconstrained interactions, surpassing any
previous study in similar conditions. These findings reveal that subordinate
fish emit chirps, which are correlated with the cessation of aggressive
behaviors. This suggests a communicative function beyond active sensing.

The introduction of this deep learning pipeline represents a significant
advancement in the methodology for studying electrocommunication, as this fully
automated approach enables the detailed analysis of large datasets and the
study of freely interacting animals in both laboratory and field settings. This
marks a step towards quantitative ethology of electrocommunication.

This following chapter will unfold in two segments: Initially, I will discuss
the implications and limitations of the behavioral correlates with the chirps
of \textit{A.~leptorhynchus}. The focus will then shift to the chirp detection
pipeline itself—evaluating its efficacy, possibilities for enhancement, and
applications in research.

\section{Are chirps a signal of submission?}

Chirps are produced in a variety of behavioral contexts. They can occur
spontaneously without any obvious external stimuli
\parencite{englerSpontaneousModulationsElectric2000a}. They can also be induced
by playing back a sine wave that mimics the electric organ discharge (EOD) of a
conspecific. Although the chirp rate decreases if the simulated conspecific
also produces chirps \parencite{walzNeuroethologyElectrocommunicationHow2013}.
Studies both in field and laboratory settings have also documented
echo-responses to chirping
\parencite{henningerStatisticsNaturalCommunication2018,
hupeEffectDifferenceFrequency2008, gamasalgadoEchoResponseChirping2011}. In the
context of mating, chirps have been observed both in the laboratory and the
field, suggesting a role in courtship and synchronization of spawning.
\textcite{hagedornCourtSparkElectric1985a} reported that a playback of male
chirps could trigger spawning in females, even without the presence of a real
male. \textcite{henningerStatisticsNaturalCommunication2018} demonstrated
stereotyped chirping patterns during courtship for two species, one recorded in
the laboratory and one in the field. They also provided observations of
chirps produced by retreating intruders that disturbed courting pairs.
Competitive interactions recorded in controlled conditions suggest opposing
functions of chirps. They have been correlated with a reduction in attack
rates, implying they might serve as a submissive signal
\parencite{hupeElectrocommunicationSpeciesWeakly2012}. In another cases,
dominant individuals were observed to produce more chirps, suggesting a role in
signaling dominance or aggression
\parencite{triefenbachChangesSignallingAgonistic2008b}. This inconsistency
along with the lack of clear behavioral correlates led
\textcite{obotiWhyBrownGhost2023} to propose that chirps might serve as probing
signals, a form of autocommunication to enhance sensory information during
social interactions.

In line with the findings of
\textcite{hupeElectrocommunicationSpeciesWeakly2012,
henningerStatisticsNaturalCommunication2018}, I have shown that in the context
of shelter competitions in \textit{A. leptorhynchus}, chirps are
disproportionally produced by fish that lost the competitions. For a signal to
be classified as communication, it needs to be a deliberate action, structure,
or trait produced by one animal (the sender) to convey information to another
animal (the receiver), influencing the receiver's behavior in a way that
typically benefits the sender
\parencite{bradburyPrinciplesAnimalCommunication2011, bergstromEvolution2016,
alcockAnimalBehavior2019}. As the chirps produced by the subordinate fish
strongly correlate with the offset of chasing events, similar to observations
by \textcite{henningerStatisticsNaturalCommunication2018}, we hypothesized that
chirps might be a means of communicating submission, and thus stopping the
chasing of the dominant. But in this case, chirps should not have correlated
precisely \emph{with} the offset of chasing, but with a point in time that is
shortly \emph{before} the offset. Additionally, ritualized chasing events
lasted longer when the loser fish produced chirps (although we cannot rule out
that the chasing events, in which chirps were produced, could have been longer
if no chirps were emitted). Because most of these chirps were produced
precisely at (or around) the end of chasing, it is sensible to assume that
prolonged chasing caused the subordinate to emit chirps, instead of the other
way around. Another explanation could be that chirps might not alter the
dominants behavior in the current competition, but increase the time it takes
for the dominant to attack again. But chirps produced by the submissive during
chasing showed no effect on the latency between attacks by the dominant.
Whether chirps of the submissive have an effect on the probability of getting
bit by the dominant has yet to be explored. Following the previous definition
of a communication signal, the reports of echo responses strongly suggest that
chirps serve as such because they change the behavior of the receiver
\parencite{henningerStatisticsNaturalCommunication2018,
hupeEffectDifferenceFrequency2008, gamasalgadoEchoResponseChirping2011}. But in
the context of the competitions analyzed here no such change in the behavior
of the receiver could be found. So while we observed clear correlations between
chirping and social status as well as with a specific phases of a ritualized
competition, it remains unclear what function chirps might serve in this
context.

The lack of evidence on the function of chirps as communication signals in this
setting does not automatically support the assumption that they do not serve a
communicative function at all. The correlations not only with social status but
also with specific spatio-temporal interactions stand in contrast with the
interpretation of \textcite{obotiWhyBrownGhost2023}, which argue that chirps
might rather be a means of active sensing through autocommunication that
enhance the perception of conspecifics. In the experiment of
\textcite{raabElectrocommunicationSignalsIndicate2021}, fish competed in
absolute darkness, leaving mainly the lateral line and electrolocation as
reliable sensory inputs. During chasing, fish stay close to each other for
prolonged times and under high velocities. The chasing fish and the one being
chased require temporally and spatially precise sensory inputs to coordinate
this movement. If chirps were emitted to perceive a conspecific, there is no
obvious explanation to why only the subordinate individual would require to do
so, specifically at the time where chasing ends. At the same time, the option
that chirps might explicitly transmit information to conspecifics does not rule
out \textcite{obotiWhyBrownGhost2023} proposal that chirps could be used for
active sensing. As suggested in other systems, such as in bats and whales,
signals that are used for active sensing could be used for communicating in
different contexts \parencite{mossAdaptiveEcholocationBehavior2023}.

The largest limitation of the competition arena by
\textcite{raabElectrocommunicationSignalsIndicate2021} is arguably that it is
confined in space. Observations of competitions in the field are constrained to
a few events observed by
\textcite{henningerStatisticsNaturalCommunication2018}, where an intruder male
interrupts a courting dyad. In response, the courting male interrupts
courtship, and approaches the intruder from a distance of more than a meter, to
which the intruder reacts with chirps and retreats. Retreat was not an option
in this study, as the fish were constrained to the aquarium. This could also be
an explanation to why competitions lasted the full three hours of the active
period in this setting. \textcite{raabElectrocommunicationSignalsIndicate2021}
showed that similar to chirps so-called frequency rises, another type of
electrocommunication signal evolving on a time scale of seconds, are
predominantly produced by subordinates and strongly correlate with the
\emph{onset} of competitions. This was interpreted in the way, that rises are
used by the subordinate to continue conspecific assessment by motivating the
dominant to compete. But while this temporal correlation of rises and chasing
onsets is strong, only 10\% of the chasing events were actually preceded by a
rise emitted by a subordinate. So the majority of competition events were
probably not initialized by electrocommunication signals but either by spatial
interactions or just spontaneously motivated by the sheer presence of another
fish in the perceivable range of the dominant individual. In the field,
territories of single fish could be larger than the experimental tank and a
scenario, in which fish are forced to compete for prolonged time, would not
occur. Following this assumption, chirps in combination with a retreat could
still signal submission in naturalistic conditions. The missing physical
possibility of retreat in the dataset could explain a lack of an effect of the
subordinate chirps onto the dominants behavior. Such multimodalities in
communication signals are extensively demonstrated for courtship signaling in
many systems and may also apply to competitions for shelter
\parencite{ronaldMaleMiceAdjust2020,
ronaldWhatMakesMultimodal2017,uetzComplexSignalsComparative2017,
redingContextdependentPreferencesVary2017,
mitoyenEvolutionFunctionMultimodal2019}. To further explore this possibility,
we either need to go back to the field to record competition scenarios in
naturalistic settings or stage competitions with the possibility of retreat.

\section{Automated detection of communication signals}

Deep learning tools for analyzing animal communication signals have been
demonstrated for the vocalizations of a wide variety of taxa including
primates, whales, rodents, bats, birds, frogs, and fish
\parencite{coffeyDeepSqueakDeepLearningbased2019,
cohenAutomatedAnnotationBirdsong2022, bellafkirBatEcholocationCall2022,
berglerORCASPOTAutomaticKiller2019, lawsonAutomatedAcousticDetection2023a,
kimuraEvaluationDeepLearningBased2023,
waddellApplyingArtificialIntelligence2021, erbsAutomatedLongtermAcoustic2023}.
In all previous cases, the models are trained entirely on tediously
hand-labeled data. The problems they solve can usually be classified into two
categories: The first one are species classifiers used on soundscape
recordings from the field, typically trained to classify spectrogram snippets
of animal vocalizations into different species during \acrfull{PAM}. These
systems often employ \acrshortpl{CNN} \parencite{kahlBirdNETDeepLearning2021,
kimuraEvaluationDeepLearningBased2023,
waddellApplyingArtificialIntelligence2021} or transformers
\parencite{bellafkirBatEcholocationCall2022}. The second category encompasses
song detectors, usually applied to some focal species in laboratory
experiments. Here, the focus is on locating the vocalizations in time (and
frequency) rather than identifying the species. Detection can be addressed
using standard CNNs \parencite{berglerORCASPOTAutomaticKiller2019}, but more
complex tasks, such as call segmentation and classification, often utilize RCNN
or YOLO models \parencite{coffeyDeepSqueakDeepLearningbased2019,
katherDevelopmentMachineLearning2024, lawsonAutomatedAcousticDetection2023a}.
An exception is TweetyNet, which incorporates \acrfull{LSTM} layers
\parencite{cohenAutomatedAnnotationBirdsong2022}.

All systems mentioned are limited to \emph{detection}. In electric fish, the
produced continuous electric field serves as an electric tag, enabling the
tracking of individuals over time \parencite{hagedornCourtSparkElectric1985a,
henningerTrackingActivityPatterns2020,
madhavHighresolutionBehavioralMapping2018,
raabAdvancesNoninvasiveTracking2022}. Since fish can be separated by frequency
and chirps are constrained to a specific relative frequency range, merely
detecting the time of chirps is insufficient. So I chose a model that predicts
chirp boundaries on the frequency axis, as previously demonstrated for rodents
\parencite{coffeyDeepSqueakDeepLearningbased2019}. Associating detected signals
with the emitting individual is a problem of \emph{assignment}. Because most
other system lack a robust passive tag, this problem remains unaddressed in any
of the presented applications. This thesis presents a first approach to emitter
assignment using feature engineering and a simple CNN for electric fish.
Together, the detection and assignment model form a continuous pipeline for
automatic chirp detection and assignment in large datasets. Instead of relying
on large hand-labeled training datasets, I additionally developed a simulation
routine to generate artificial recordings of electric fish. This way, I can
automatically generate training data from a small representative sample of
communication signals. This reduces the need for the error-prone process of
hand-labeling large amounts of training data.

\subsection{Performance of the detection and simulation pipeline}

In terms of performance, both the metrics used to evaluate the detection model
as well as the assignment step are promising: The average precision reaches
more that 85\% (AP@[.5:.01:.95], figure \ref{fig:detector_performance}) for
detection and 97\% for assignment (figure \ref{fig:assignment_training}). From
these metrics, detection and assignment errors are expected to be quite rare.
But it is important to emphasize that the detection performance is evaluated
partially on simulated data and the assignment model fully on simulated data,
which probably leads to overestimating the performance of the pipeline. For
this reason, I manually checked hundreds of spectrograms in addition to the
performance metrics, two of which are introduced in detail in figures
\ref{fig:chirp_detections_simultaneous} and
\ref{fig:chirp_detections_bigchirps}. These spectrograms illustrate that for
most cases and in recordings with an adequate \acrfullpl{SNR}, the detection
pipeline performs much better compared to previous implementations,
particularly for longer and larger chirps. First applications on field data
should paint a clearer picture on how usable this approach is and what needs to
be improved. A promising approach to further increase the performance of the
detection model could be through active learning. Here, the model labels
previously unseen data. From these detections, uncertain samples, such as those
with low confidence scores, are picked, manually annotated and the model is
retrained on those. This is a thoroughly tested method of increasing the
performance of machine learning algorithms
\parencite{zhanComparativeSurveyDeep2022}.

In addition to a thorough evaluation of the detection pipeline it is equally
important to evaluate and improve the simulation pipeline. Electric fish offer
the unique opportunity to capture their full behavior by electric recordings
but also to simulate a complete scene of multiple interacting fish. In such
scenarios, even a trained human observer struggles to assign fish to the
correct emitter. This makes the availability of \emph{true} ground-truths that
do not rely on human annotators indispensable. The current simulation routine
is merely a first proof of principle, where many aspects of natural electrode
grid recordings are still missing.

Currently, it starts by generating EOD$f$ traces that include communication
signals such as chirps and rises. Based on those frequency traces, the EOD is
generated from Fourier components to resemble that of \textit{A.
leptorhynchus}. A separate pathway generates movement of this simulated fish
and the EOD is then amplitude-modulated on each simulated electrode to reflect
the current distance between the fish and each electrode.

To improve this pipeline, we can start at the simulation of communication
signals. Chirps are currently generated by extracting the evolution of the
EOD$f$ from chirps of real fish and inserting them into the simulated EOD$f$ of
the simulated fish. This way, the diversity of the simulations of chirps is
constrained by the available data. Instead of using the extracted instantaneous
frequencies of chirps directly, we could train a learning
algorithm to generate artificial EOD$f$s of chirps. Algorithms such as
\acrfullpl{VAE} \parencite{nishikimiVariationalAutoencoderBased2024} or
\acrfullpl{GAN} \parencite{dissanayakeGeneralizedGenerativeDeep2023,
hazraSynSigGANGenerativeAdversarial2020} are already used in medical
applications to synthesize EEG, EMG and ECG signals and could be applied in a
similar fashion to synthesize the frequency evolution of the EOD during chirps
or rises. Instead of the real EOD$f$ traces, we could then use the synthetic
communication signals to simulate fish.

The second step would be the improvement of the simulation on the level of the
EOD itself. In the current simulation, the fish is an oscillating monopole and
the electric field decays equally in all directions. But in reality, the
electric field is an oscillating dipole with two lobes, which does not decay
equally in all directions \parencite{bendaPhysicsElectrosensoryWorlds2020}.
This difference would not be visible on the sum spectrograms I used for chirp
detection, which is why I ignored it for now. But when future versions of chirp
detectors incorporate spatial information into the assignment step, this
feature could become important. Multiple models of such an electric field
already exist and could be implemented and incorporated into the current
pipeline \parencite{babineauModelingElectricField2006,
nelsonPreyCaptureWeakly1999, wallachInternalModelCanceling2023}. As a result,
the way in which the field decays would not just depend on position, but on
body size and curvature of the fish.

After adding realistic simulation of the dipole, we should also generate the
fish movement based on the statistics of real movements. This is
important, because these motions substantially influence the amplitude
modulations of the EOD. These are visible on spectrograms and perceivable by
other fish \parencite{yuElectrosensoryContrastSignals2019,
fotowatStatisticsElectrosensoryInput2013a}. In the current simulation, only the
distribution of swimming velocity is taken from approximations from field
recordings. Other patterns, such as forward-backward motions and heading
directions, are hand-tuned only to \emph{resemble} those of a real fish.
Instead, we should use data from experiments where animals were video-tracked.
This would allow us to generate models that not only produce realistic
velocities but also trajectories, body curvature, and maybe even behavioral
patterns such as foraging, fighting or resting. All of those would bring us
closer to reliably recreate the full statistics of natural recordings and hence
closer to true ground truths to build and benchmark analysis algorithms such as
the chirp detector.

One of the most important aspects of the present simulation pipeline is the
augmentation with real noise. Adding realistic noise to simulated recordings is
not only essential to enable the transfer to real data but also provides
endless possibilities to \emph{augment} the simulations. The same simulation on
many different \enquote{noise backdrops} can drastically increase the number of
training samples for learning algorithms
\parencite{salamonScaperLibrarySoundscape2017}. Currently, this noise is taken
from another recording of fish in the laboratory. Instead, we should aim to
create a library of noise backdrops from different settings in the lab an in
the field. This would not require much effort, since the recording hardware
already exists and fish would not be needed. Not just the background noise
itself but also the \acrfull{SNR} is crucial. The current simulation pipeline
is restricted to a single SNR. Instead, we could use various levels of SNRs to
generate many training samples for a single simulation-noise combination, which
would make the model more robust to weak EOD amplitudes or fish that are far
away from the electrodes.

The last step would be to incorporate the complexity of the environment in such
simulations. Many objects substantially distort the electric field
\parencite{zlenkoVisualizationElectricFields2023,
pereiraContextualEffectsSmall2005, bacherNewMethodSimulation1983} and the
natural habitat of \textit{A. leptorhyncus} is a complex three dimensional
labyrinth of stacked rocks \parencite{hagedornCourtSparkElectric1985a}. The
surface structure of such habitats could be reconstructed using underwater
photogrammetry \parencite{franciscoHighresolutionNoninvasiveAnimal2020}. The
problem with this approach is that fish prefer to hide in crevices between
those rocks, which are usually not picked up well by conventional
photogrammetry approaches. How such a complex 3D environment could be included
in these simulations is yet to be explored.

Improving the simulation pipeline would also facilitate a comparison between
the chirp detector's and human-level performance, the gold-standard in
supervised detection and classification tasks. This way of evaluating
performance is hardly ever viable, because in most cases, the ground truth is
generated by human annotators. This way, a model logically never achieves to
surpass human level performance. But in a case where the actual ground truth is
accessible, such as in simulations, human level performance can be compared to
the detectors performance. Under noisy conditions, this could enable the
training of detection pipelines that could surpass the performance of human
annotators.

\subsection{Convolutional neural networks versus transformer architectures}

RCNN and YOLO architectures, such as the ones used in this pipeline, have
already been used successfully for the automatic detection of ultrasonic
vocalizations of rodents \parencite{coffeyDeepSqueakDeepLearningbased2019}.
Here, I provide another successful application of object detection models to
the detection of electric signals. But in the last years, the development of
CNNs for computer vision and audio analysis has slowed down due to the rise of
transformer-based architectures. Transformers were mainly developed for natural
language processing and became popular through the public demonstration of the
\acrfull{GPT}. Since then, transformer-based architectures have been
successfully applied to vision tasks and usually outperform convolution-only
architectures in terms of accuracy and precision
\parencite{shehzadiObjectDetectionTransformers2023}. But convolutional
approaches still surpass transformers, when the computational time is
considered. In the work presented here, the processing time had to be as low as
possible to facilitate the fast processing of weeks of recordings. But
transformers will likely become more performant in the future. The
computational costs will decrease and the available hardware will become
better. At that point, we should move to transformer models for
Neuroethological signal recognition.

\subsection{Back to the time domain}

Wave-type weakly electric fish offer the unique property that in opposite to
bats or whales, the signals they produce are continuously emitted. Furthermore,
each EOD approximately consists of just a single frequency at a point in time.
As mentioned in \hyperref[fig:multifish_chirps]{chapter two}, if we record a
single fish, we can estimate this frequency for a very fine temporal precision
by the inverse of each period. But as soon as noise is introduced or the
signals of multiple fish overlap, this becomes impossible as illustrated in
figure \ref{fig:multifish_chirps}. So
\textcite{raabAdvancesNoninvasiveTracking2022,
henningerTrackingActivityPatterns2020} developed the \texttt{wavetracker} to
approximate the frequency changes over time of multiple fish on spectrograms. I
followed with the \texttt{chirpdetector} that is capable of also detecting
chirps on spectrograms of multiple fish. But working on these spectrograms
comes with a large pitfall: We are always restricted by the time-frequency
uncertainty. As an example, the spectrograms used in the \texttt{chirpdetector}
have a temporal resolution of \SI{20}{\milli\second}, corresponding to a single
frequency estimate for short chirps. This is just enough to make them visible
but still reduces the frequency resolution in a way that fish less than
\SI{20}{\hertz} apart are hardly separable from each other. But what if we were
able to extract each fish from such a recording similarly to how we can record
single fish in the laboratory? I believe that this is not so far away as we
might think: Convolution and transformer architectures running on spectrograms,
waveforms or both simultaneously can be used for sound source separation,
meaning that they can separate e.g., vocals, bass and drums from a rock song
\parencite{rouardHybridTransformersMusic2022,
stollerWaveUNetMultiScaleNeural2018, petermannCocktailForkProblem2022}.
Applications of sound source separation in Biology have already been
demonstrated successfully for vocalizations of birds and fish
\parencite{dentonImprovingBirdClassification2022,
mancusiFishSoundsEvaluation2022}. Our problem is a bit different. In the case
of separating instruments from vocals (or bird songs of different species)  the
task is \emph{semantic} segmentation. If we would try to apply this approach to
electric fish, i.e., separating a recording of multiple fish into recordings of
single fish, we would be faced with \emph{instance} segmentation. Instead of
segmenting EODs from the background, the task is to segment the EODs of
individual fish from each other \emph{and} the background. While this sounds
complex, the fact that fish are usually separated in frequency could make this
approachable. But we would need a learning algorithm that combines methods from
sound source separation with methods from instance segmentation from computer
vision. An improved simulation pipeline would be an ideal way of training such
models. Having temporal resolutions not constrained to that of spectrograms
would be an enormous benefit that allows for much more intricate insights into
the social lives of electric fish. This is because we do not just use frequency
to understand their communications, but also amplitude to estimate their
positions.

\subsection{Applications of chirp detection}

With these technical considerations in mind, for what other practical
applications could the \texttt{chirpdetector} be useful? Electric fish are one
of the main fractions that make up the total biomass of fish in the Amazon
basin \parencite{lassoCompositionAbundanceBiomass1992,
cramptonEcologicalPerspectiveDiversity2011a}. They present a large variability
in habitat use, social lives and communication repertoires that are
comparatively poorly explored and still offer exiting discoveries to be made.
With the right training data, the chirp detector could be trained and used on
the chirps of any chirping species. With the simulation pipeline to
generate training data, just a small representative subset might be enough to
get good results. In combination with an electrode grid and the
\texttt{wavetracker}, we now have the tools to record and analyze electric and
spatio-temporal behaviors of multiple, freely interacting wave-type weakly
electric fish in their natural habitat.
\textcite{henningerStatisticsNaturalCommunication2018} demonstrated that this
approach is indispensable in reflecting upon decades of neuroethological
research in electric fish. \textcite{fortuneSpookyInteractionDistance2020a}
showed that this approach can be employed to study behavioral differences
between surface and cave dwelling populations of two closely
related species. With the pipeline presented here, we can add the quantitative
analysis of chirps to this suite of methods. Just by the sheer amount of chirps
that are produced in various behavioral contexts, without arguing for or
against a certain function, it is sensible to assume that they are
important to the lives of electric fish
\parencite{henningerStatisticsNaturalCommunication2018, obotiWhyBrownGhost2023,
hupeElectrocommunicationSignalsFree2009}. Because chirps are connected
tightly to courtship or spawning, they could be used as proxies for mating or
at least courtship events
\parencite{henningerStatisticsNaturalCommunication2018,
hagedornCourtSparkElectric1985a}. Currently not even the mating system of the
different species of wave-type electric fish is known. Beyond ethological and
comparative field studies, this approach is equally applicable to recordings of
multiple fish in controlled conditions. \textcite{obotiWhyBrownGhost2023}
correctly criticized the lack of causal evidence for the function of chirps as
a communication signal. The pipeline presented here should make it much easier
to study chirps and enable a much higher throughput compared to manual
annotations, to shed a light onto the function(s) of chirps.

In parallel to fundamental research, the chirps of weakly electric fish could
become equally important in the applied sciences, such as in population
monitoring. The electric sense is widely regarded as a significant evolutionary
boon, contributing substantially to the remarkable diversity and abundance of
electric fish \parencite{cramptonEvolutionElectricSignal2006}. This unique
adaptation, however, necessitates specific conditions in terms of nutritional
availability and water chemistry, rendering these species potentially
vulnerable to anthropogenic disturbances, as discussed by
\textcite{markhamEnergeticsSensingCommunication2016}. An aspect not covered in
their discussion is the potential exacerbation of these vulnerabilities by
electrical noise. In areas close to human settlements, recordings of electric
fish in their natural habitats are often marred by significant electrical
interference within their perceptible frequency ranges—a problem not
encountered in more remote locations. A particularly concerning source of such
interference is the proliferation of hydroelectric dams
\parencite{agostinhoAllThatGoes2011, keppelerEarlyImpactsLargest2022,
duponchelleConservationMigratoryFishes2021}. These structures not only pose
barriers to riverine migration but also generate substantial electrical noise.
Monitoring fish populations is essential to prevent further damage to their
diversity and abundance \parencite{larentisEffectsHumanDisturbance2022}. Active
monitoring by telemetry is possible but expensive, error prone and restricted
to few large species \parencite{hahnBiotelemetryRevealsMigratory2019}.
Monitoring through eDNA, where species accounts are extracted from the DNA
available in water samples, is a promising and holistic approach and removes
the need for more invasive monitoring methods
\parencite{desantanaCriticalRoleNatural2021}. To complement these approaches,
purely passive monitoring was demonstrated recently in the endangered Amazonian
River Dolphins, using hydrophones and neural networks to detect echolocation
signals \parencite{erbsAutomatedLongtermAcoustic2023}. But this is restricted
to two species. I would argue that we can take this a step further with
wave-type, and maybe at some point even pulse-type electric fish.
\textcite{henningerStatisticsNaturalCommunication2018} and
\textcite{madhavHighresolutionBehavioralMapping2018} have demonstrated that
electric fish can not only be passively monitored, but that we can even track
individual identities over time and estimate their positions to a comparatively
fine temporal scale. What has been missing were their communication signals. If
we would be able to monitor and automatically analyze the communication of
these fish, we would not only be able to asses how their abundances change over
time but also how the reproductive behavior of these populations is affected by
environmental changes. This could be implemented in a purely passive way,
without the need to catch fish or significantly disturb their environment.
 
\section{Conclusion} % (fold)
\label{sec:conclusion}

The automated chirp detection pipeline developed in this thesis successfully
detects all known types of chirps in recordings of multiple Brown Ghost
Knifefish. This represents a significant advancement in the study of
electrocommunication, since it is the first system capable of detecting and
automatically assigning each chirp to the fish that emitted it.

Applying this pipeline to a dataset from staged competitions between two fish
surpassed the number of chirps in most previous studies and revealed important
behavioral insights. Chirps were predominantly emitted by subordinate fish and
were correlated with the cessation of aggressive interactions by dominant
individuals. This finding suggests a communicative function of chirps in
mediating ritualized competitions for resources.

In conclusion, the chirp detection pipeline lays the groundwork for deeper
exploration into the electrocommunication of wave-type electric fish. It paves
the way for studying social dynamics, behavioral cues, and ecological roles of
chirping in both laboratory and natural settings and marks a first step towards
non-invasive and passive monitoring of electrocommunication in wild
populations.

% section Conclusion (end)

% chapter Discussion (end)

\iffalse
The disparity between the current capabilities in data analyses and the
practices with which data is not only analyzed but also recorded in some fields
of research is striking. While controlled experiments and
\enquote{conventional} statistics are indispensable, whether such causal
effects revealed in experiments play a role in the natural environment and
behavioral context of species is often overlooked in the push for causality in
strictly controlled conditions. But especially now, as not only the autonomous
recording of enormous amounts of data but also automated analyses become more
and more accessible, going into the field and inspiring new hypotheses by
quantitative observations becomes approachable. As an example, smartphone apps
can automatically recognize bird songs and classify the species just from
comparatively poor-quality phone recordings
\parencite{kahlBirdNETDeepLearning2021}. Despite this, many studies still rely
on manually labeling communication signals of all kinds of species. Software
such as
\href{http://www.mackenziemathislab.org/deeplabcut}{\texttt{DeepLabCut}} and
\href{https://sleap.ai/}{\texttt{SLEAP}} enables automated tracking of
individual body parts of most animals \parencite{pereiraSLEAPDeepLearning2022a,
lauerMultianimalPoseEstimation2022}. This provides the possibility of high
throughput in behavioral experiments and frees time that can be allocated to,
for example, recording more data. But computational technologies are still not
widely adopted nor included into teaching curricula. Instead, many rely on
cheap manual labor in the form of undergraduate students to annotate datasets
by hand. \textcite{andersonScienceComputationalEthology2014} nicely summarized
this process as beeing slow, imprecise, subjective, low-dimensional, limited by
the human sensory system and lastly \emph{dull}. The need for annotated
datasets is not something that we soon be able to completely resolve, but
instead of labeling a full dataset by hand, it is much less time consuming to
label a small representative subset to train a machine learning system that
then labels the full dataset for analysis. This would decouple the amount of
data that can be recorded from the number (or patience) of available
annotators. This way, junior scientists, often times still students themselves,
would not be prevented from learning skills that become more and more valuable
not only in science but already are so in industry.

\fi


% This ensures that the subsequent sections are being included as root
% items in the bookmark structure of your PDF reader.
\bookmarksetup{startatroot}
\backmatter

\begingroup
\let\clearpage\relax
\glsaddall
\printglossary[type=\acronymtype]
% \newpage
% \printglossary
\endgroup

\printindex
\printbibliography

\chapter{Appendix} % (fold)
\label{appendix}

\begin{center}
  \textsc{Extracting chirps for model fitting in full simulations}
\end{center}

I used the following pipeline to extract estimates on how the instantaneous
frequency evolves during a chirp using the output of the sliding convolutional
classifier as a prior on the Raab dataset:

\begin{enumerate}

  \item Electrode selection: By the power of the frequency band of the EOD of
    the fish I extracted the signal on the electrode of the strongest power,
    i.e., the electrode that the fish was closest to. For each chirp, I took a
    full second, meaning 500 \si{\milli\second} before and after the chirp. The
    result is a single voltage trace over one second, i.e. 20,000 samples, as
    the signals were sampled at \SI{20}{\kilo\hertz}.

  \item Band-pass filtering: I then applied a band-pass filter after
    multiplying the snipped with a Tukey window ($\alpha=0.1$) to reduce edge
    effects of the filter. The cutoff frequency below the EOD frequency
    of the upper fish EOD$f_{\text{up}}$ was set relative to that of the lower fish EOD$f_{\text{low}}$ at $\text{EOD}f_{\text{low}} + 0.8 \times
    \Delta\text{EOD}f$ and above the EOD $f_{\text{up}}$ at $1.9
    \times$EOD$f_{\text{low}}$. The lower cutoff frequency was chosen to reduce
    interference from the EOD of the other fish while not cutting off potential
    down-modulations of the frequency that sometimes occur after chirps. The
    upper cutoff was chosen to reduce the interference from the first harmonic
    of the EOD of the lower fish while not cutting off the top of the chirp.

  \item Standardization: I then standardized the filtered EOD to have a mean of zero
    and a standard deviation of one.

  \item Instantaneous frequency: I then computed the instantaneous frequency of
    the signal using zero crossings. The instantaneous frequency was
    additionally smoothed using a Gaussian kernel with a standard deviation of
    5 samples. It was then linearly interpolated to the original sample rate of
    the signal. The result is an estimate on how the EOD$f_{up}(t)$ evolved
    over the course of one second EOD$f_{\text{inst}}(t)$ at the temporal
    resolution of a single period of the EOD of the upper fish.

  \item Centering baseline frequency: To center the baseline of this snippet of
    the instantaneous frequency that contains a chirp to zero, I estimated the
    baseline frequency from the largest peak of the distribution of all
    frequencies in the snippet. To estimate it, I used a histogram of 100 bins
    as this was much faster than using a kernel density estimate. The baseline
    frequency estimate was then subtracted from the instantaneous frequency to
    shift the baseline frequency to zero.

  \item Peak detection: The most descriminating feature of chirps in
    Apteronotids is the rapid increase of the EOD$f$ during a chirp. To
    estimate the exact time of the chirp, I detected peaks in the instantaneous
    frequency that exceeded the 99\th percentile of the centered instantaneous
    frequencies using the \texttt{find\_peaks} function of the \texttt{scipy}
    package \parencite{virtanenSciPyFundamentalAlgorithms2020}. Before
    detecting the peaks, I applied another Tukey window ($\alpha=0.3$) to the
    centered instantaneous frequencies to scale down fluctuations at the edges
    of the snippet (as chirps are expected in the center of the snippet). I
    only considered the largest peak and discarded the whole snippet to
    continue with the next chirp when the peak was more than 250
    \si{\milli\second} away from the center of the snippet. I also discarded
    the whole snippet if included multiple peaks that where larger than $0.3
    \times \max{\text{EOD}f_{\text{snip}}}$.

  \item Centering chirp time: To estimate a center of the peak that is
    independent of the shape of the peak, I descended both the left and right
    flanks until I reached $0.1 \times \max{\text{EOD}f_{\text{snip}}}$. The
    position of the chirp was then defined as the center of the two points on
    the flanks. The snippet was then rolled to center the chirp position within
    the snippet. As the edges could now be non-zero again, I applied another
    Tukey window ($\alpha=0.3$) to scale down fluctuations at the edges of the
    snippet. The result is a zero-centered estimate of how the frequency of the
    EOD of the fish evolves over one chirp with EOD-period resolution.

\end{enumerate}

This was repeated for each chirp for six of the recordings in the competition
dataset, in which the $\Delta\text{EOD}f$ was sufficiently large and resulted
in a total of 615 chirps. Each snippet was linearly interpolated to the
original sampling rate of \SI{20}{\kilo\hertz}.

\newpage

\begin{center}
  \textsc{Extracting chirps for hybrid simulations}
\end{center}

The following steps were taken to extract the instantaneous frequency of the
chirps for every recording in the chirp-chamber dataset and parts of the Oboti
dataset:

\begin{enumerate}

  \item Signal preprocessing: To find the EOD$f$ of the fish I used peak
    detection on the power spectrum between 300 and 1200 \si{\hertz}. The
    signal was then band-pass filtered between $\text{EOD}f - 100\si{\hertz}$
    and $\text{EOD}f + 1.9\text{EOD}f$. to remove all harmonics of the EOD.

  \item Instantaneous frequency and envelope: To estimate the instantaneous
    frequency of the EOD, I used the zero crossings of the signal to get the
    start and stop times of each period of the EOD. The instantaneous frequency
    was then computed as the inverse of the difference of the start and stop
    times of each period. The envelope of the EOD was estimated by taking the
    maximum value of each period of the EOD. Both the instantaneous frequency
    and the envelope were then linearly interpolated to the original sample
    rate of the signal.

  \item Smoothing and detection: I smoothed the instantaneous frequency to
    different degrees depending on the task: To detect peaks in the
    instantaneous frequency as chirps, I used a Gaussian kernel with a standard
    deviation of 30 samples. Peaks with a prominence between 20 and 400
    \si{\hertz} as well as a width between 4 and 4000 samples were then
    detected using the \texttt{find\_peaks} function of the \texttt{scipy}
    package \parencite{virtanenSciPyFundamentalAlgorithms2020}. I then
    extracted the snippet of the centered instantaneous frequency 200 ms before
    and after the peak. The result is a 2D matrix were each row is a snippet of
    the smoothed instantaneous frequency that contains a chirp. The same
    applies to the envelopes. Meanwhile, I always also kept the original
    instantaneous frequency and envelope.

  \item Shifting and normalization: To make all instantaneous frequencies
    relative to the EOD$f$ of the emitter, I shifted them to zero by
    subtracting the first element of each row. I then normalized the
    envelope by dividing it by the maximum value of each row. In the following,
    the detected chirps passed a series of filters to remove false positives.

  \item Filtering:
    The following filters were applied to the detected chirps:
    \begin{enumerate}

      \item Median envelope: I removed all chirps that had a median envelope
        lower than 0.2. Low median amplitudes indicate that the amplitude
        modulation in the window contains a peak instead of a trough.

      \item Frequency peak: I removed all chirps that had a peak frequency of
        more than 550 \si{\hertz}.

      \item Frequency trough: While an undershoot of the frequency does often
        happen after a chirp, it does not usually go further than 100
        \si{\hertz} below the baseline frequency. I removed all chirps that
        had a trough frequency of less than 200 \si{\hertz}.

      \item Baseline frequency: To look at how stable the baseline is in the
        chirp snippet, I looked at the first and last 100 \si{\milli\second}
        of each chirp, which should be the baseline frequency and mostly zero.
        If the 5th percentile of this part of the snippet was lower than
        \SI{-30}{\hertz}, I removed the chirp.

      \item Baseline envelope: If the same window on the normalized envelope
        had standard deviation of more than 0.2 I removed the chirp, as this
        would correspond to a very noisy signal.

      \item Envelope trough: I removed all chirps that had a trough of the
        envelope of less than 0.1 because the lower the envelope of the signal,
        the less precise the instantaneous frequency estimate. I encountered
        many cases where peaks in the instantaneous frequency where just
        artifacts that appeared when the amplitude of the signal approached
        zero.

      \item Chirp width: If the width of the chirp was less than 9
        \si{\milli\second}, I removed the chirp.

      \item If the width to height ratio was less than
        \SI{0.001}{\second\per\hertz}, I removed the chirp.

    \end{enumerate}

  \item Centering time: To find the center of the chirp, I computed a Boolean
    array that was true for all values above 20 \si{\hertz} and false for all
    values below 20 \si{\hertz}. If the number of peaks in this array was more
    than 1, or if the peak reached up to one edge of the snippet, I removed
    the chirp. I then determined the center of the chirp as the center between
    the two threshold crossings. To center the snippet, I then rolled the
    snippet so that the center of the chirp was at the center of the snippet.

\end{enumerate}

\end{document}
